% concatenated from paper.tex, tables/summary_nums.tex, tables/pertype.tex
%% Classification of compact hyperbolic Coxeter 6-polytopes with 10 facets.
%% Compile:  pdflatex paper && bibtex paper && pdflatex paper && pdflatex paper
\documentclass[11pt,reqno]{amsart}

\usepackage[T1]{fontenc}
\usepackage[utf8]{inputenc}
\usepackage{amsmath,amssymb,amsthm,mathtools}
\usepackage{booktabs}
\usepackage{longtable}
\usepackage{tikz}
\usepackage{array}
\usepackage{microtype}
\usepackage[margin=1.15in]{geometry}
\usepackage[colorlinks=true,linkcolor=blue!55!black,citecolor=blue!55!black,
            urlcolor=blue!55!black]{hyperref}

\usetikzlibrary{calc}

\theoremstyle{plain}
\newtheorem{theorem}{Theorem}[section]
\newtheorem{proposition}[theorem]{Proposition}
\newtheorem{lemma}[theorem]{Lemma}
\newtheorem{corollary}[theorem]{Corollary}
\theoremstyle{definition}
\newtheorem{definition}[theorem]{Definition}
\newtheorem{remark}[theorem]{Remark}
\newtheorem{assumption}[theorem]{Assumption}

\DeclareMathOperator{\conv}{conv}
\DeclareMathOperator{\rank}{rank}
\DeclareMathOperator{\sig}{sign}
\newcommand{\HH}{\mathbb{H}}
\newcommand{\RR}{\mathbb{R}}
\newcommand{\QQ}{\mathbb{Q}}
\newcommand{\ZZ}{\mathbb{Z}}
\newcommand{\Gr}{G}
\newcommand{\PP}{P_{6,10}}

\newcommand{\NumOrderTypes}{14{,}309{,}547}
\newcommand{\NumTypes}{387}
\newcommand{\NumTypesMF}{381}
\newcommand{\NumSurvivors}{54}
\newcommand{\NumSearched}{52}
\newcommand{\NumSearchedLess}{51}
\newcommand{\NumDegKilled}{76}
\newcommand{\NumBigMF}{7}
\newcommand{\NumPrism}{119}
\newcommand{\NumAfterComb}{185}
\newcommand{\NumAfterDegMF}{304}
\newcommand{\NumRequired}{11}
\newcommand{\NumRedundant}{43}
\newcommand{\NumRedundantSearched}{41}
\newcommand{\DFiveWild}{141{,}033}
\newcommand{\DFiveWildKilled}{141{,}024}
\newcommand{\DFiveWildSurvived}{9}
\newcommand{\WildCertified}{406}
\newcommand{\WildLabelOnly}{387}
\newcommand{\WildRootBox}{19}
\newcommand{\WildBnB}{0}
\newcommand{\WildMinMargin}{1.1\times10^{-3}}
\newcommand{\WildResidMin}{9.1\times10^{-5}}
\newcommand{\WildResidMax}{3.3\times10^{-4}}
\newcommand{\DFourCascade}{822{,}432}
\newcommand{\DFourPair}{132{,}571}
\newcommand{\DFourFallback}{132{,}571}
\newcommand{\PlainRejected}{545}
\newcommand{\PlainCertified}{545}
\newcommand{\DFiveWildCertSet}{2{,}411}
\newcommand{\DFiveWildCert}{2{,}400}
\newcommand{\DFiveWildUndecided}{11}
\newcommand{\DFiveWildScreenPass}{11}
\newcommand{\DFiveWildPassRefuted}{0}
\newcommand{\DFiveWildCensus}{141{,}033}
\newcommand{\DFiveWildJointKilled}{2{,}400}
\newcommand{\NumCascade}{546}
\newcommand{\NumPlain}{546}
\newcommand{\NumWild}{406}
\newcommand{\NumRequiredLess}{10}
\newcommand{\NumThreeFree}{2}
\newcommand{\ThreeFreeList}{159, 329}
\newcommand{\NumSubtrees}{11}
\newcommand{\NumAssignments}{952}
\newcommand{\CPUHours}{0.96}
\newcommand{\RealizingType}{379}
\newcommand{\RealizingP}{3}
\newcommand{\RealizingProfile}{2^{3}3^{3}5^{3}}
\newcommand{\ThreeFreeP}{6, 6}
\newcommand{\ThreeFreeProfile}{2^{6}}
\newcommand{\NumGramConfigs}{12}
\newcommand{\AutTypeOrder}{24}
\newcommand{\AutStabOrder}{2}
\newcommand{\WeightLeaves}{1}
\newcommand{\DFiveTotal}{51}
\newcommand{\DFiveTypes}{109}
\newcommand{\DFiveCPUHours}{2.4}
\newcommand{\MaxDepth}{0}
\newcommand{\DFiveRealizing}{0\!\to\!22, 6\!\to\!18, 19\!\to\!6, 29\!\to\!3, 5\!\to\!1, 63\!\to\!1}
\newcommand{\DFourFound}{348}
\newcommand{\DFourTarget}{348}


%% ---- Coxeter-diagram drawing conventions (matching Burcroff, Ma--Zheng) ----
%% single line = pi/3, double = pi/4, triple = pi/5, dashed = ultraparallel.
\newcommand{\cnode}[2]{\fill (#1) circle (2.4pt) node[#2] {}}
\newcommand{\esingle}[2]{\draw[line width=.5pt] (#1)--(#2);}
\newcommand{\edouble}[2]{%
  \draw[line width=1.9pt] (#1)--(#2);%
  \draw[line width=0.8pt,white] (#1)--(#2);}
\newcommand{\etriple}[2]{%
  \draw[line width=2.9pt] (#1)--(#2);%
  \draw[line width=1.9pt,white] (#1)--(#2);%
  \draw[line width=1.0pt] (#1)--(#2);%
  \draw[line width=0.35pt,white] (#1)--(#2);}
\newcommand{\edash}[2]{\draw[line width=.5pt,dashed] (#1)--(#2);}

\begin{document}

\title[Compact hyperbolic Coxeter six-dimensional polytopes with ten facets]
      {Compact hyperbolic Coxeter six-dimensional polytopes\\ with ten facets}

\author{John Mcleod}
\address{Solvo.ai}
\email{john@solvo.ai}

\date{\today}

\subjclass[2020]{Primary 51F15; Secondary 20F55, 52B11, 52B40}
\keywords{hyperbolic Coxeter polytope, reflection group, Gale diagram,
          order type, computer-assisted classification}

\begin{abstract}
We show that, up to isometry, there is exactly one compact hyperbolic Coxeter
$6$-polytope with $10$ facets: the polytope $\PP$ whose Coxeter diagram appears
as Figure~5 of Burcroff \cite{Burcroff2024}, attributed there to Bugaenko
\cite{Bugaenko1984}. Together with the results
of Lann\'er, Kaplinskaja and Esselmann, Felikson--Tumarkin (no such polytopes for
$d\ge 8$; uniqueness for $d=7$), and the independent classifications of
Burcroff and Ma--Zheng in dimensions $4$ and $5$ ($348$ and $51$ polytopes),
this completes the classification of compact hyperbolic Coxeter $d$-polytopes
with $d+4$ facets in every dimension. The proof is computer assisted. We
enumerate all $\NumTypes$ candidate combinatorial types from the complete
database of planar order types on $10$ points via affine Gale duality, reduce to
$\NumRequired$ types using two combinatorial consequences of Lann\'er's
classification, the classification of compact Coxeter $d$-polytopes with $d+2$
facets, and the classification of compact Coxeter $5$-polytopes
with $9$ facets, and decide those $\NumRequired$ by an exhaustive search over
Coxeter
labellings with no a~priori bound on the dihedral angles. The search terminates
with machine-checked exhaustion certificates covering every searched type; the unique surviving Gram matrix is certified exactly over
$\QQ(\sqrt2,\sqrt5)$ and independently by \texttt{CoxIter}. The emptiness verdicts
are exact as well, and in two stages: forward checking, which empties ten of the
$\NumRequired$ types on its own, decides ellipticity and Lann\'er membership
combinatorially and by integer and $\QQ(\sqrt2,\sqrt3,\sqrt5)$ arithmetic; and every
labelling that reaches a screen and is not accepted is refuted exactly over that
field --- for all but $\WildRootBox$ of them by the non-vanishing of a single
determinant, and otherwise by interval arithmetic on a compactified domain. So no
verdict rests on a floating-point tolerance, on a numerical search, or on a bound on
the ultraparallel weights. The same code path,
unchanged, reproduces both completed censuses in this family exactly: the $51$
polytopes in dimension $5$ and the $348$ in dimension $4$.
Code, data, per-type certificates and re-run instructions are publicly
available. Most of the software was written by an AI assistant under the
author's direction; see \S\ref{sec:provenance}.
\end{abstract}

\maketitle

\section{Introduction}
\label{sec:intro}

Let $\HH^d$ denote real hyperbolic $d$-space. A \emph{hyperbolic Coxeter
polytope} is a convex polytope $P\subset\HH^d$, bounded by finitely many
geodesic hyperplanes, all of whose dihedral angles have the form $\pi/m$ with
$m\in\ZZ_{\ge 2}$. Such a polytope is the fundamental domain of the discrete
group generated by reflections in its facets. We are concerned with
\emph{compact} such polytopes, and we write $n$ for the number of facets and
$k = n-d$.

The classification of compact hyperbolic Coxeter $d$-polytopes with $d+k$ facets
has been completed one value of $k$ at a time. For $k=1$ (simplices) this is
Lann\'er \cite{Lanner1950}; for $k=2$ it is Kaplinskaja \cite{Kaplinskaja1974}
and Esselmann \cite{Esselmann1996}; for $k=3$ it is Esselmann
\cite{Esselmann1994} and Tumarkin \cite{Tumarkin2004}. For $k=4$, Felikson and
Tumarkin \cite{FeliksonTumarkin2008} proved that no such polytope exists for
$d\ge 9$ (their Theorem~1) nor for $d=8$ (their Theorem~2), and that the
$7$-dimensional example with $11$ facets, found by Bugaenko
\cite{Bugaenko1984}, is unique (their Theorem~3). Burcroff \cite{Burcroff2024}
and, independently and by different methods, Ma and Zheng
\cite{MaZheng2024,MaZheng2023} classified the cases $d=4$ and $d=5$, finding
$348$ and $51$ polytopes respectively. This left exactly one gap. In Burcroff's
words \cite[\S1]{Burcroff2024}, ``the only remaining dimension where new
polytopes may arise is $d=6$'', where ``the only known polytope was constructed
by Bugaenko''.

Our main result closes that gap.

\begin{theorem}
\label{thm:main}
Up to isometry there is exactly one compact hyperbolic Coxeter $6$-polytope with
$10$ facets, namely the polytope $\PP$ whose Coxeter diagram is shown in
Figure~5 of \cite{Burcroff2024}.
\end{theorem}

Theorem~\ref{thm:main} is proved by an exhaustive computer search whose
reductions are justified in \S\ref{sec:completeness}; the assumptions on which
it rests are collected and discussed candidly in \S\ref{sec:assumptions}.
Combining it with the results quoted above gives the following.

\begin{corollary}
\label{cor:complete}
The classification of compact hyperbolic Coxeter $d$-polytopes with $d+4$ facets
is complete in every dimension. Such polytopes exist precisely for
$2\le d\le 7$ \cite{FeliksonTumarkin2008}; in dimensions $2$ and $3$ there are
infinitely many, and for $d\ge4$ their numbers up to isometry are
\[
\begin{array}{c|ccccc}
d      & 4   & 5  & 6 & 7 & \ge 8\\\hline
\#     & 348 & 51 & \mathbf{1} & 1 & 0
\end{array}
\]
\end{corollary}

The polytope $\PP$ is not new: it is drawn in \cite[Fig.~5]{Burcroff2021} and
again in \cite[Fig.~5]{Burcroff2024}, in both places attributed to Bugaenko
\cite{Bugaenko1984}, and we verified independently that the
polytope produced by our search is isomorphic to that figure edge for edge and
label for label.
What is new here is that it is the \emph{only} one, and that the search space in
which this is proved is explicitly delimited.


\subsection*{Method, in outline}
The proof has the shape common to \cite{Burcroff2024,MaZheng2024,MaZheng2023}:
first enumerate the possible combinatorial types, then decide, for each type,
which Coxeter labellings yield a Gram matrix of signature $(6,1)$.

A compact Coxeter polytope is simple (Proposition~\ref{prop:simple}), so a
$6$-polytope with $10$ facets has a $3$-dimensional Gale diagram and hence a
planar affine Gale diagram. By Burcroff's Theorem~3.5, every such polytope
arises from a planar order type on $10$ points with two interior points marked
positive. The order types on $10$ points have been completely enumerated: there
are $\NumOrderTypes$ of them \cite{AichholzerAurenhammerKrasser2002,
AichholzerKrasser2001,OrderTypeDB}. Processing all of them with an
exact-arithmetic affine Gale criterion yields $\NumTypes$ combinatorial types
(\S\ref{sec:types}).

The decisive reduction is a facet filter (\S\ref{sec:filters}). Every facet of
a compact Coxeter $6$-polytope is a compact Coxeter $5$-polytope, and Ma and
Zheng proved that exactly six simple $5$-polytopes with $9$ facets admit such a
structure \cite[\S6]{MaZheng2023}. Applying this combinatorially cuts
$\NumTypes$ types to $\NumRequired$, of which one realizes a polytope. Of the
$\NumSurvivors$ passing the facet filter alone, Lemmas~\ref{lem:degree},
\ref{lem:mfsix} and~\ref{lem:prism} remove $\NumRedundant$; what remains is decided
by search. Among the $\NumRedundant$ are the $\NumThreeFree$ types all of whose
missing faces have size $2$, which Esselmann's bound (Lemma~\ref{lem:esselmann})
would also exclude; that bound is therefore not needed here, and we record it only
because it is the cleanest reason no such type can occur.

The search (\S\ref{sec:search}) enumerates Coxeter labellings with
forward-checking the elliptic and Lann\'er conditions of
Proposition~\ref{prop:faces}, using
Ma--Zheng's wildcard device \cite[Prop.~3.5]{MaZheng2024} so that
\emph{no a~priori bound is placed on the dihedral angles}; edges of label
$\ge 7$ are carried as a continuous unknown and resolved by an integer window
scan. Acceptance is exact: $100$-digit Gauss--Newton refinement of the
ultraparallel weights, exact signature and local-isolation certification, and
finally an independent \texttt{CoxIter} check. Rejection is exact too. Forward
checking decides its two predicates without a tolerance
(Proposition~\ref{prop:exactgates}), and it is what empties ten of the
$\NumRequired$ types: only the realizing type produces a labelling that reaches a
screen at all. Each of the $\NumAssignments$ that do is refuted exactly
(Propositions~\ref{prop:plainexact} and~\ref{prop:wildexact}): substituting
$x_e=1/t_e$ compactifies the
ultraparallel directions, each $(d+2)$-principal minor becomes a polynomial with
coefficients in $\QQ(\sqrt2,\sqrt3,\sqrt5)$ on a closed rational box containing
the domain, and outward-rounded interval arithmetic certifies that one such minor
vanishes nowhere on it --- for the $\PlainRejected$ labellings without a wildcard
the certifying minor involves no unknown at all, so the refutation is the
non-vanishing of a single determinant over that field. The search consumed
$\CPUHours$ CPU-hours, screened $\NumAssignments$ labellings, and emitted an
exhaustion certificate for each of the $\NumRequired$ types. Exactly one type, and
within it exactly one polytope, survived.

\section{Preliminaries}
\label{sec:prelim}

We follow the conventions of \cite{Vinberg1985,Burcroff2024}.

\subsection{Gram matrices and diagrams}
Let $P=\bigcap_{i\in I} f_i^-\subset\HH^d$ be a Coxeter polytope with facets
$f_i$ and outward unit normals $e_i$ in the Lorentzian space $\RR^{d,1}$. The
\emph{Gram matrix} of $P$ is $\Gr = (\langle e_i,e_j\rangle)_{i,j\in I}$, so
\[
\Gr_{ii}=1,\qquad
\Gr_{ij}=
\begin{cases}
-\cos(\pi/m_{ij}) & \text{if $f_i,f_j$ meet at angle $\pi/m_{ij}$},\\
-1                & \text{if $f_i,f_j$ are parallel},\\
-\cosh\rho_{ij}<-1& \text{if $f_i,f_j$ diverge at distance $\rho_{ij}$.}
\end{cases}
\]
The associated \emph{Coxeter diagram} $\Sigma(P)$ has a node for each facet;
nodes $i,j$ are joined by an edge labelled $m_{ij}$ when $f_i\cap f_j\neq
\emptyset$ and $m_{ij}\ge 3$, left unjoined when $m_{ij}=2$, and joined by a
\emph{dashed} edge when the facets diverge. Following \cite{Burcroff2024} we
draw an edge with $m=3,4,5$ as a single, double, triple line respectively, and
label larger $m$ numerically. An edge with $m\le 5$ is said to have \emph{low
weight}; one with $m\ge 6$ is \emph{multi-multiple}.

A subdiagram is \emph{elliptic} if its Gram matrix is positive definite,
\emph{parabolic} if it is positive semidefinite and degenerate, \emph{Lann\'er}
if it is connected, all of its proper subdiagrams are elliptic, and its Gram matrix
is indefinite --- equivalently $\det\Gr<0$, which is what excludes the parabolic
case --- and \emph{superhyperbolic} if its Gram matrix has negative inertia
index $>1$. The last clause is not redundant: the affine diagram
$\smash{\widetilde{G_2}}$, a path with labels $3$ and $6$, is connected, is not
elliptic and has every proper subdiagram elliptic, yet its Gram determinant
vanishes, so it is parabolic and not Lann\'er. Dropping the clause would make
Corollary~\ref{cor:mfsize} false, since $\smash{\widetilde{A_6}}$, a $7$-cycle of
edges labelled $3$, would then count as a Lann\'er diagram of order $7$.

\begin{proposition}[Vinberg; see {\cite[Remark~4.3]{Burcroff2024}}]
\label{prop:vinberg}
If $\Sigma=\Sigma(P)$ is the Coxeter diagram of a compact hyperbolic Coxeter
$d$-polytope $P$ with $n$ facets, then $\Sigma$ is connected, contains no
parabolic subdiagram and no bold edge, is not superhyperbolic, and $\Gr(\Sigma)$
has signature $(d,1)$ and rank $d+1$; equivalently, its kernel has dimension
$n-d-1$.
\end{proposition}

For $d=6$, $n=10$ this reads: $\Gr$ is a $10\times10$ symmetric matrix of
signature $(6,1)$, rank $7$, kernel dimension $3$.

\subsection{Faces and missing faces}
A subset $J\subseteq I$ is a \emph{face} of $P$ if $\bigcap_{i\in J}f_i\neq
\emptyset$; a \emph{missing face} is a minimal non-face.

\begin{proposition}[{Vinberg; \cite[Prop.~4.4]{Burcroff2024}}]
\label{prop:faces}
Faces of codimension $m$ in $P$ correspond to elliptic subdiagrams of order $m$
in $\Sigma(P)$, and the missing faces of $P$ correspond to the Lann\'er
subdiagrams of $\Sigma(P)$.
\end{proposition}

\begin{proposition}[{Vinberg; \cite[Remark~4.5]{Burcroff2024}}]
\label{prop:simple}
Every compact hyperbolic Coxeter $d$-polytope is simple: every $(d-j)$-face is
contained in precisely $j$ facets. In particular every vertex lies on exactly
$d$ facets.
\end{proposition}

This is immediate from Proposition~\ref{prop:faces}, since an elliptic
subdiagram of order $d$ has exactly $d$ nodes.

Since Lann\'er diagrams exist only in orders $2\le \ell\le 5$
\cite{Lanner1950}, Proposition~\ref{prop:faces} gives at once:

\begin{corollary}
\label{cor:mfsize}
Every missing face of a compact hyperbolic Coxeter polytope has size between $2$
and $5$. Missing faces of size $2$ are exactly the pairs of disjoint facets,
i.e.\ the dashed edges of $\Sigma(P)$.
\end{corollary}

We write $p=p(P)$ for the number of pairs of disjoint facets.

\section{Combinatorial types via order types and Gale diagrams}
\label{sec:types}

\subsection{The order-type route}
By Proposition~\ref{prop:simple} a compact Coxeter $6$-polytope with $10$ facets
is a simple $6$-polytope with $10$ facets; dually, a simplicial $6$-polytope
with $10$ vertices. Its Gale diagram is a configuration of $10$ points in
$\RR^{3}$, which can be reduced to a planar \emph{affine} Gale diagram: a
configuration of $10$ labelled points in $\RR^2$, each marked positive or
negative \cite{Gale1956,Ziegler1995}. The relevant statement is:

\begin{theorem}[{\cite[Thm.~3.5]{Burcroff2024}}]
\label{thm:burcroff35}
Every compact Coxeter $d$-polytope $P$ with $d+4$ facets admits an affine Gale
diagram obtained by taking an arrangement of $d+4$ points $A\subseteq\RR^2$ in
general position and choosing two points from the interior of $\conv(A)$ to be
positive. Moreover, the combinatorial type of $P$ is completely determined by
the order type of $A$.
\end{theorem}

Thus enumerating combinatorial types reduces to enumerating order types of $10$
points, together with the $\binom{10}{2}$ choices of positive pair.

\subsection{Completeness of the order-type database}
\label{sec:aak}

The \emph{order type} of a planar point set in general position is its
chirotope, i.e.\ the orientation of every ordered triple. Order types realizable
by actual point sets have been completely enumerated for $n\le 11$ by
Aichholzer, Aurenhammer and Krasser
\cite{AichholzerAurenhammerKrasser2002,AichholzerKrasser2001}; the resulting
Order Type Data Base \cite{OrderTypeDB} stores one integer-coordinate realizing
point set per order type. For $n=10$ it contains $\NumOrderTypes$ order types.
This is the count quoted by Burcroff \cite[\S3]{Burcroff2024} and is the
database she uses for $d=4,5$; we use the same files
(\texttt{otypes10.b16}). The enumeration is exhaustive for realizable order
types by construction: it proceeds by inserting points one at a time into all
realizable order types on $n-1$ points and keeping one representative per
equivalence class, and every stored representative comes with explicit integer
coordinates, so every entry is realizable and no realizable order type is
omitted. Theorem~\ref{thm:burcroff35} produces the affine Gale diagram of an
actual polytope, whose underlying point configuration is realizable; hence the
database covers every combinatorial type we must consider.

\subsection{From order types to combinatorial types}
For each order type and each candidate positive pair $\{u,v\}$ we test, using
exact integer arithmetic on the stored coordinates, that both $u$ and $v$ lie in
the interior of the convex hull of the remaining eight points, and then read off
the face lattice by the standard affine Gale criterion \cite[\S6.4]{Ziegler1995}:
$S\subseteq[n]$ is a face if and only if, in $T=[n]\setminus S$, the positive
points of $T$ and the negative points of $T$ are ``interleaved'' in the sense
that the relative interiors of $\conv(T^+)$ and $\conv(T^-)$ intersect. This
criterion is implemented over the rationals; no floating-point comparison enters the
determination of a face. Types are then deduplicated exactly: two types are
identified when a permutation of $[10]$ carries one missing-face hypergraph to
the other, tested by colour refinement followed by an exhaustive search over
colour-preserving permutations. This is legitimate as a test of combinatorial
equivalence because a simplicial complex is determined by its minimal non-faces:
$S$ is a face if and only if $S$ contains no missing face, so the missing-face
hypergraph determines the whole face lattice, and hence the combinatorial type of
the dual simple polytope.

That argument uses the \emph{full} hypergraph, and the generator records minimal
non-faces only up to size $5$ (\S\ref{sec:completeness}), so deduplication is
carried out on truncated data and we owe an argument that nothing was lost. Two
distinct types with isomorphic $\le5$ hypergraphs, differing only in a minimal
non-face of size $6$, would be merged; that would matter if the surviving
representative were then discarded by Lemma~\ref{lem:mfsix}, which excludes types
possessing a minimal non-face of size $6$, since a merged partner without one
would be discarded with it. It does not happen. For each of the $\NumBigMF$ types
that Lemma~\ref{lem:mfsix} excludes, we recover every source candidate that was
merged into it --- $2{,}124$ of them in total --- and verify that each one
\emph{also} has a minimal non-face of size $6$. So every candidate the lemma
removes deserves removal on its own data, and no type is lost to the truncation.

Two properties of the input are checked rather than assumed. The stored
coordinates are verified to be in general position---no three of the ten points
collinear, i.e.\ no vanishing triple determinant---using exact integer arithmetic;
and for every generated type the vertex list is compared against the $6$-subsets of
facets containing no \emph{recorded} missing face. Because minimal non-faces are
recorded only up to size $5$ (see \S\ref{sec:completeness}), the two need not
coincide: a $6$-subset avoiding every recorded missing face may fail to be a vertex,
which is exactly the signature of a minimal non-face of size $6$ and is how
Lemma~\ref{lem:mfsix} is applied. What the comparison does verify unconditionally is
the other direction---every vertex has size $6$ and contains no missing face---which
is the consistency of the two representations.

\begin{proposition}
\label{prop:387}
This procedure yields exactly $\NumTypes$ combinatorial types of simple
$6$-polytopes with $10$ facets and at least two pairs of disjoint facets.
\end{proposition}

The number $p$ is distributed over the $\NumTypes$ types as
$p=2$: 245, $p=3$: 101, $p=4$: 28, $p=5$: 10, $p=6$: 3.

\begin{remark}
Burcroff established the order-type route in all dimensions and carried out the
enumeration for $d=4$ and $d=5$, and in dimension $6$ analysed the case in which
the missing faces have sizes in $\{2,5\}$ \cite[\S3, \S8]{Burcroff2024}. Of our
$\NumTypes$ types, $\NumTypesMF$ contain a missing face of size $3$ or $4$.
\end{remark}

\subsection{Validation of the generator}
\label{sec:genvalid}

Run on $n=8$ and $n=9$, the same generator reproduces the published combinatorial
censuses exactly: $30$ types for $d=4$ and $109$ for $d=5$, matching the
ground-truth lists up to relabelling, with no spurious types. Both checks are
exact, by permutation search on missing-face hypergraphs. These are the
lists of \cite{GruenbaumSreedharan1967} and \cite{MaZheng2023} respectively.

\begin{remark}[three corrections to {\cite[Appendix~A]{Burcroff2024}}]
\label{rem:appendixA}
The second table of that appendix, the $d=5$ candidate list, has $111$ rows, which
describe $109$ combinatorial types. Two rows carry identical missing-face lists;
two further rows are not the missing-face systems of simple polytopes; and one type
present in our output and in \cite{MaZheng2023} has no row. After those three
corrections her list and ours agree exactly. None of them affects her
classification, and each is documented with its data in \cite{Repo}. (The first
table of the same appendix is the $34$-row $d=4$ list used in
\S\ref{sec:d4valid}.)
\end{remark}

\section{Restrictions on combinatorial types and Coxeter diagrams}
\label{sec:filters}
\label{sec:completeness}

This section justifies every reduction applied between the $\NumTypes$ types of
Proposition~\ref{prop:387} and the exhaustive search of \S\ref{sec:search}. We
state each as a lemma, with a proof or a precise citation. A reduction is
\emph{sound} if it never discards a type that admits a compact Coxeter
structure.

The first condition is standard and is used throughout \cite{Burcroff2024}; we
record the derivation only because it draws on two separate results and we know
of no source stating it in this combined form.

\begin{lemma}[at least two disjoint pairs]
\label{lem:p2}
A compact hyperbolic Coxeter $d$-polytope with $d+4$ facets has $p\ge 2$.
\end{lemma}
\begin{proof}
A compact Coxeter polytope with $p=0$ is a simplex or one of the seven
Esselmann $4$-polytopes with $6$ facets
\cite[Thm.~A]{FeliksonTumarkin2008}; in either case $n\le d+2$. By the Main
Theorem of \cite{FeliksonTumarkin2009}, quoted as
\cite[Thm.~5.1]{Burcroff2024}, a compact hyperbolic Coxeter $d$-polytope with
exactly one pair of non-intersecting facets has at most $d+3$ facets. Hence
$n=d+4$ forces $p\ge2$.
\end{proof}

This condition is imposed inside the generator, which is why
Proposition~\ref{prop:387} already counts only types with $p\ge2$. It is the
only combinatorial condition applied there; the following are applied
afterwards.

By Corollary~\ref{cor:mfsize} every missing face has size in $\{2,3,4,5\}$. All
$\NumTypes$ generated types already satisfy this, so it eliminates nothing; we
recall it because it is what makes the ``size $3$ or $4$'' condition below
equivalent to a statement about $3$-free polytopes.

Two purely combinatorial exclusions come next. Both are theorems, both are cheap,
and between them they remove more than a quarter of the types that would otherwise
have to be searched.

\begin{lemma}[bounded dashed degree]
\label{lem:degree}
Let $P$ be a compact hyperbolic Coxeter $d$-polytope with $n$ facets, $d\ge6$, and
let $f_i$ be a facet disjoint from exactly $t$ others. Then $t\le n-d-2$. For
$d=6$, $n=10$ this says no facet is disjoint from more than two others: the dashed
edges of $\Sigma(P)$ form a graph of maximum degree at most $2$.
\end{lemma}
\begin{proof}
$P$ is simple, so for each of the $n-1-t$ facets $f_j$ meeting $f_i$ the
intersection $f_i\cap f_j$ is a facet of $f_i$, and distinct $j$ give distinct
facets; hence $f_i$ is a compact hyperbolic Coxeter $(d-1)$-polytope with exactly
$n-1-t$ facets. A compact hyperbolic Coxeter polytope of dimension $d-1\ge5$ is not
a simplex, since by Lann\'er's classification \cite{Lanner1950} compact hyperbolic
Coxeter simplices exist only in dimensions $2$, $3$ and $4$. So $f_i$ has at least
$d+1$ facets, giving $n-1-t\ge d+1$.
\end{proof}

\begin{lemma}[no missing face of size $d$]
\label{lem:mfsix}
No compact hyperbolic Coxeter $6$-polytope has a missing face of size $6$.
\end{lemma}
\begin{proof}
By Proposition~\ref{prop:faces} a missing face of size $6$ would correspond to a
Lann\'er subdiagram of order $6$, and by \cite{Lanner1950} Lann\'er diagrams exist
only in orders $2\le\ell\le5$.
\end{proof}

\begin{lemma}[a facet disjoint from two others is a prism]
\label{lem:prism}
Let $P$ be a compact hyperbolic Coxeter $6$-polytope with $10$ facets and let $f_i$
be disjoint from exactly two other facets. Then $f_i$ is combinatorially the
simplicial prism $\Delta^4\times I$, and in particular its missing-face profile is
$2^1 5^1$.
\end{lemma}
\begin{proof}
As in Lemma~\ref{lem:degree}, $f_i$ is a compact hyperbolic Coxeter $5$-polytope
with $9-2=7$ facets, that is with $d+2$ facets for $d=5$. Compact hyperbolic
Coxeter $d$-polytopes with $d+2$ facets are completely classified: they are the
simplicial prisms of \cite{Kaplinskaja1974} together with the products of two
simplices of \cite{Esselmann1996}, and the latter occur only for $d=4$. Hence
$f_i\cong\Delta^4\times I$. Its facets are the five ``sides'' $F\times I$, one for
each facet $F$ of $\Delta^4$, and the two ``ends'' $\Delta^4\times\{0\},
\Delta^4\times\{1\}$; the ends are disjoint, giving a missing face of size $2$, and
the five sides have empty common intersection while any four of them meet, giving a
minimal non-face of size $5$. There are no others, so the profile is $2^1 5^1$.
\end{proof}

This is a strong condition, because the profile of $f_i$ is determined by the
combinatorial type of $P$ (Lemma~\ref{lem:facetmf}) and can simply be read off. Of the
$\NumTypes$ types it eliminates $\NumPrism$ beyond those already removed by
Lemmas~\ref{lem:degree} and~\ref{lem:mfsix}.

The second lemma is not vacuous, and it is worth being explicit about why. The
generator records the minimal non-faces of each type up to size $5$; a type may
therefore possess a minimal non-face of size $6$ without one appearing in its
recorded list, and this is detected by comparing the vertex list against the
$6$-subsets of facets containing no recorded missing face. Of the $\NumTypes$
types, $\NumDegKilled$ are excluded by Lemma~\ref{lem:degree} and a further
$\NumBigMF$ by Lemma~\ref{lem:mfsix}, leaving $\NumAfterDegMF$; Lemma~\ref{lem:prism}
then removes $\NumPrism$ more, leaving $\NumAfterComb$.

\begin{lemma}[facet admissibility]
\label{lem:facet}
Let $P$ be a compact Coxeter $6$-polytope with $10$ facets and let $f_i$ be a
facet meeting all nine others. Then the combinatorial type of $f_i$, a simple
$5$-polytope with $9$ facets, has one of the four missing-face size profiles
\[
2^6 5^3,\qquad 2^5 3\,4\,5^2,\qquad 2^4 3^2 4^2 5,\qquad 2^3 3^3 4^3 .
\]
\end{lemma}
\begin{proof}
Each facet of a compact Coxeter polytope is itself a compact Coxeter polytope of
one lower dimension (the reflection group of $P$ restricted to the hyperplane of
$f_i$; see \cite[Prop.~4.6]{Burcroff2024} and \cite{Vinberg1985}). If $f_i$
meets all nine other facets, then, $P$ being simple, each $f_i\cap f_j$ is a
facet of $f_i$ and distinct $j$ give distinct facets, so $f_i$ has exactly nine
facets and is a
compact Coxeter $5$-polytope with $9$ facets. Ma and Zheng prove that ``only six
of all simple $5$-polytopes with $9$ facets admit compact hyperbolic structure''
\cite[\S6, Tables~13--14]{MaZheng2023}; these are their $P_{312}$, $P_{319}$,
$P_{322}$, $P_{302}$, $P_{313}$, $P_{284}$. We computed the missing-face
profiles of these six directly from their combinatorial types and obtained
exactly the four profiles listed (the first is shared by $P_{312},P_{319},
P_{322}$, the second by $P_{302}$, the third by $P_{313}$, the fourth by
$P_{284}$).
\end{proof}

Lemma~\ref{lem:facet} is a genuine weakening of ``$f_i$ is one of the six'': it
tests only the profile, so it retains any type whose facet merely shares a
profile with a realizing $5$-polytope. It therefore errs strictly on the side of
keeping types.

The profile of $f_i$ must be computed from the incidence structure and not read off
the missing faces of $P$, because the naive rule is wrong: if $\{i,a,b\}$ is a
missing face of $P$ then $\{a,b\}$ is a missing face of $f_i$, yet $\{a,b\}$ is a
\emph{face} of $P$. The correct rule is the following, and it is what the
implementation uses.

\begin{lemma}[missing faces of a facet]
\label{lem:facetmf}
Let $P$ be simple and let $f_i$ be any facet. Identify the facets of $f_i$ with
$N(i)=\{j\neq i: \{i,j\}\text{ is a face of }P\}$ via $f_j\mapsto f_i\cap f_j$.
Then $J\subseteq N(i)$ is a face of $f_i$ if and only if $J\cup\{i\}$ is a face of
$P$; consequently the missing faces of $f_i$ are the minimal non-faces of the
link of $i$ in the simplicial complex dual to $P$. In particular, when $f_i$ meets
every other facet, $N(i)=[n]\setminus\{i\}$.
\end{lemma}
\begin{proof}
$\bigcap_{j\in J}(f_i\cap f_j)=f_i\cap\bigcap_{j\in J}f_j$, which is non-empty
exactly when $J\cup\{i\}$ is a face of $P$; and $f_i\cap f_j=\emptyset$ for
$j\notin N(i)$, so those $j$ index no facet of $f_i$. Minimal non-faces of a
complex are determined by its faces, which gives the second statement.
\end{proof}

Applied to all facets meeting all nine others---detected
combinatorially as the facets occurring in no size-$2$ missing face---it reduces
the $\NumTypes$ types to $\NumSurvivors$, and combined with
Lemmas~\ref{lem:degree}, \ref{lem:mfsix} and~\ref{lem:prism} to $\NumRequired$. This is by far the
largest single reduction, and it is exactly as strong as Ma--Zheng's $d=5$
classification; see \S\ref{sec:assumptions}\ref{as:facet}.

We record the resulting cascade, since it is what delimits the search:
\[
\begin{array}{rcl}
\NumTypes &\text{types}& \text{generated (Proposition~\ref{prop:387})}\\
-\,\NumDegKilled && \text{Lemma~\ref{lem:degree}, bounded dashed degree}\\
-\,\NumBigMF && \text{Lemma~\ref{lem:mfsix}, no missing face of size }6\\
-\,\NumPrism && \text{Lemma~\ref{lem:prism}, degree-2 facets are prisms}\\
\NumAfterComb && \text{remain}\\
\NumRequired && \text{also pass the facet filter, and are searched.}
\end{array}
\]
Every step is a theorem, and none of them touches the realizing type.

That last point is worth more than an assurance, and it can be measured.
$\NumRedundantSearched$ of the $\NumRedundant$ excluded types have also been put
through the search directly, along with the $\NumRequired$ that survive: all
$\NumSearched$ are rigorous zeros by exhaustion except the realizing type, which
yields its one polytope. So the exclusions discard nothing realizable, and the
evidence is independent of their proofs. (The $\NumThreeFree$ not covered are the
$3$-free ones. They are excluded by Lemma~\ref{lem:degree} like the rest, but they
carry no Lann\'er constraint at all, so their label trees do not terminate in
reasonable time and the search cannot corroborate them --- which is why a theorem is
wanted there.) Table~\ref{tab:pertype} reports every
survivor with the lemma that decides it; the search data for the excluded types are
in \cite{Repo}.

\begin{lemma}[Esselmann; $3$-free polytopes need $2d$ facets]
\label{lem:esselmann}
A compact Coxeter $d$-polytope all of whose missing faces have size $2$ (a
\emph{$3$-free} polytope) has at least $2d$ facets, with equality only for the
$d$-cube.
\end{lemma}
\begin{proof}
This is \cite[Lemma~6.7]{Esselmann1994}, quoted as
\cite[Lemma~10.6]{Burcroff2024}.
\end{proof}

\begin{corollary}
\label{cor:threefree}
No compact Coxeter $6$-polytope with $10$ facets is $3$-free.
\end{corollary}
\begin{proof}
A $3$-free compact Coxeter $6$-polytope needs at least $12$ facets, and $10<12$.
\end{proof}

Exactly $\NumThreeFree$ of the $\NumSurvivors$ types passing the facet filter are
$3$-free, both with $p=6$ and missing-face profile $\ThreeFreeProfile$. Each of
them is already excluded by Lemma~\ref{lem:degree}, so Corollary~\ref{cor:threefree}
is not needed for the classification; we retain it because it is the cleanest
statement of why no $3$-free type can occur here at all.

\begin{remark}
Burcroff proves \cite[Thm.~8.1]{Burcroff2024} that there is no compact Coxeter
$6$-polytope with $10$ facets whose missing faces have orders only $2$ and $5$,
and states in her abstract the consequence that such a polytope must contain a
missing face of size $3$ or $4$. Her proof of Theorem~8.1 proceeds by analysing
the two combinatorial types she exhibits there, both of which have a missing
face of size $5$; it does not literally cover the $3$-free case, in which no
missing face of size $5$ occurs. (We do not reuse her labels for these two types:
the symbol $I_2$ is needed here for the dihedral Coxeter group $I_2(m)$.) We therefore invoke Lemma~\ref{lem:esselmann}
instead, which settles the $3$-free case unconditionally. Among our
$\NumSurvivors$ surviving types, the only ones without a missing face of size
$3$ or $4$ are the $\NumThreeFree$ $3$-free ones, so
Corollary~\ref{cor:threefree} supersedes any appeal to
\cite[Thm.~8.1]{Burcroff2024}, which is consequently not used in this paper.
\end{remark}

\begin{lemma}[low-weight edges]
\label{lem:burcroff55b}
Let $v_1v_2$ be an ordinary edge of an admissible abstract Coxeter diagram
$\Sigma$. If $\Sigma$ has no Lann\'er diagram containing both $v_1$ and $v_2$,
has a Lann\'er diagram $L$ of order greater than $2$ containing $v_2$, and has
no dashed edge from $v_1$ to any vertex of $L$, then $v_1v_2$ has low weight
$(m\le5)$.
\end{lemma}
\begin{proof}
\cite[Lemma~5.5(b)]{Burcroff2024}.
\end{proof}

By Proposition~\ref{prop:faces}, the Lann\'er subdiagrams of $\Sigma(P)$ are
exactly the missing faces of $P$, and the dashed edges are exactly the
size-$2$ missing faces. Hence the hypotheses of Lemma~\ref{lem:burcroff55b}
depend only on the combinatorial type, and the resulting caps $m\le 5$ can be
computed once per type and imposed during enumeration. This is a pruning device only; it
never removes a type.

Finally, the enumerator requires $\Sigma(P)$ to be connected; this is part of
Proposition~\ref{prop:vinberg} (\cite[Remark~4.3]{Burcroff2024}).

\subsection{The label alphabet is not bounded a~priori}
\label{sec:alphabet}

Restricting ordinary edge labels to a finite set is the most attackable step in
classifications of this kind, and we avoid it. We know of no published theorem
bounding the labels of a compact Coxeter $d$-polytope with $d+4$ facets for
$d\ge5$; Burcroff's Lemma~5.5 bounds only those edges satisfying its
hypotheses, and the alphabet $\{2,\dots,10,12\}$ observed in the finished
$d=4$ and $d=5$ censuses is an a~posteriori observation, not a theorem.
Instead we use Ma--Zheng's truncation.

\begin{proposition}[{\cite[Prop.~3.5]{MaZheng2024}}]
\label{prop:mz35}
There are finitely many finite Coxeter groups of rank $r$ with Coxeter matrix
entries at most seven. Consequently it suffices to enumerate potential matrices
with entries at most seven, the other candidates being obtained by substituting
integers greater than seven with the value seven.
\end{proposition}

Concretely, our enumeration uses the alphabet $\{2,3,4,5,6,\star\}$ in which
$\star$ is a wildcard standing for ``any $m\ge7$'', evaluated during
forward-checking at $\cos(\pi/7)$. We record the exact sense in which this is
lossless.

\begin{lemma}
\label{lem:wildcard}
Let $S$ be a set of nodes carrying labels in $\{2,\dots,6\}$ on some pairs and
labels $m_e\ge7$ on a set $W\neq\emptyset$ of pairs.
\begin{enumerate}
\item[(a)] The Gram matrix of $S$ is positive definite for some choice of
$(m_e)_{e\in W}$ with all $m_e\ge7$ if and only if it is positive definite when
every $m_e=7$; indeed the verdict does not depend on $(m_e)$ at all.
\item[(b)] If $|S|=3$, the Gram matrix of $S$ has exactly one negative
eigenvalue for some choice of $(m_e)_{e\in W}$ if and only if it does when every
$m_e=7$.
\item[(c)] If $|S|\ge4$, no Lann\'er diagram on $S$ contains an edge of label
$\ge6$. \emph{(This is part of Lann\'er's classification
\cite{Lanner1950}; see the remark after the proof.)}
\end{enumerate}
\end{lemma}
\begin{proof}
(a) By the classification of finite Coxeter groups, a connected elliptic diagram
of rank $\ge3$ has all labels $\le5$ (the connected spherical types of rank
$\ge3$ are $A_r,B_r,D_r,E_{6,7,8},F_4,H_3,H_4$, the Coxeter groups of those
names---not to be confused with Burcroff's combinatorial types $H_i$). Hence in an elliptic diagram
any edge of label $\ge6$ must be an $I_2(m)$ connected component. So $S$ is
elliptic if and only if each $e\in W$ is a connected component $I_2(m_e)$ of $S$
and the rest of $S$ is elliptic; $I_2(m)$ is elliptic for every finite $m$, so
the verdict is independent of $(m_e)$. If some $e\in W$ is not a component, $S$
is non-elliptic for every $m_e\ge6$.

(b) Write $S=\{p,q,r\}$ with $1/p+1/q+1/r$ the usual triangle-group invariant
and $r=m_e\ge7$ the wild label. If $\{p,q\}=\{2,2\}$ the diagram is
$A_1\times I_2(r)$, elliptic for every $r$, and never has a negative
eigenvalue. Otherwise $1/p+1/q\le 1/2+1/3=5/6$ and
$1/p+1/q+1/r\le 5/6+1/7<1$ for every $r\ge7$, so the triangle is hyperbolic
with exactly one negative eigenvalue, for every $r\ge7$. Either way the verdict
at $r=7$ agrees with the verdict at every $r\ge7$. (When two or three edges are
wild the same dichotomy applies, since $1/p+1/q+1/r$ only decreases.)

(c) Suppose $L$ is a Lann\'er diagram of order $k\ge4$ and $uv$ is an edge of
$L$ with label $\ge6$. For any third node $w$ of $L$, the subdiagram on
$\{u,v,w\}$ is a proper subdiagram, hence elliptic, hence by the argument in
(a) it is $A_1\times I_2(m)$, so $w$ is orthogonal to both $u$ and $v$. Since
$k\ge4$ such a $w$ exists, and as this holds for every $w\notin\{u,v\}$ the
pair $\{u,v\}$ is a connected component of $L$, contradicting the connectedness
of a Lann\'er diagram.
\end{proof}

\begin{remark}
Part (c) is not new: it is contained in Lann\'er's classification
\cite{Lanner1950}, which lists every Lann\'er diagram of order $4$ and $5$ and
in which no edge label exceeds $5$. We include the short argument only because
we rely on the statement in the precise form above, and we have in addition
checked it computationally rather than by reading the published figure:
enumerating all connected rank-$3$ diagrams
with labels up to $200$ confirms that none with an edge of label $\ge6$ is
elliptic (the boundary case being the parabolic diagram
$\smash{\widetilde{G_2}}$ with labels $3,6$, whose Gram determinant
$1-\cos^2(\pi/3)-\cos^2(\pi/6)$ vanishes identically), and enumerating from the
definition over labels $\{2,\dots,5\}$ recovers Lann\'er's lists exactly: $9$
diagrams of order $4$ and $5$ of order $5$, largest edge label $5$.
\end{remark}

Part (c) is what lets a missing face of size $\ge4$ carrying an edge of label
$\ge6$ be rejected outright: such a configuration is not the Lann\'er diagram of a polytope
whatever the label is, so the rejection is not merely harmless but exact, and
$\cos(\pi/7)$ never has to be evaluated there. Parts (a) and (b) make the
substitution at $m=7$ lossless for the remaining cases, where the label enters only
through the elliptic test and the three-node Lann\'er criterion. As an independent
check we evaluated both
predicates at $m\in\{7,8,9,10,12,15,20,30,50,100,300,1000\}$ for every label
tuple on $3$ and $4$ nodes over $\{2,\dots,6,\star\}$ with at least one wild
edge ($91$ and $31{,}031$ tuples respectively) and for $8000$ random tuples on
$5$ nodes---about two million $(\text{tuple},m)$ evaluations in all. The verdict
never differed from the verdict at $m=7$.

Wildcard-bearing labellings are not decided at $m=7$; they are resolved by the
range analysis described in \S\ref{sec:search}.

\section{The search}
\label{sec:search}

Fix a surviving combinatorial type $T$ with missing faces $M_1,\dots,M_r$. Its
dashed edges are the $M_i$ of size $2$; write $D$ for that set and $k=|D|$. The
unknowns are an integer label $m_e\ge2$ for each ordinary pair $e$ and a real
weight $x_e>1$ for each $e\in D$, and the constraint is
$\rank\Gr = 7$ with signature $(6,1)$.

\subsection{Enumeration with forward checking}
Ordinary labels are enumerated by backtracking over the alphabet
$\{2,3,4,5,6,\star\}$ in a fixed order that places the edges lying in every
missing face first. Two families of constraints are forward-checked:

\begin{itemize}
\item \emph{Elliptic groups.} For each face of $T$ of size $3$, $4$ or $6$---in
      particular for each vertex, a set of $6$ facets---the corresponding Gram
      submatrix must be positive definite (Proposition~\ref{prop:faces}). The
      condition is on \emph{faces}, not on pairwise-meeting sets: three pairwise
      intersecting facets with empty triple intersection form a missing face of
      size $3$, whose subdiagram is Lann\'er and therefore not positive definite.
\item \emph{Lann\'er groups.} For each missing face of size at least $3$ the Gram
      submatrix must be Lann\'er. (Such a missing face has no dashed pair: by
      minimality it cannot contain a missing face of size $2$.)
\end{itemize}
Only the sizes $3$, $4$ and $6$ are tested, and nothing is lost by that. Every face
of a simple $6$-polytope is contained in a vertex and positive definiteness is
inherited by principal submatrices, so the vertex condition on $6$ facets implies the
condition on faces of every smaller size; the sizes $3$ and $4$ are tested in
addition purely because they prune earlier in the enumeration order, and $1$, $2$ and
$5$ are omitted for the same reason they are redundant.

For each group the set of admissible label tuples is precomputed once and stored
as a bitmask indexed by (edge position, label), so a partial assignment is
extended in $O(1)$ integer operations per group. Edges capped by
Lemma~\ref{lem:burcroff55b} are restricted to $m\le5$ at this stage.

Both predicates are decided \emph{exactly}, with no eigenvalue tolerance anywhere,
which matters because the forward-checking prune is what decides all but one of the
$\NumRequired$ searched types.

\begin{proposition}
\label{prop:exactgates}
Ellipticity and Lann\'er membership are decidable by exact arithmetic on the
labelled diagram, for every label including the wildcard, and the search decides
them that way.
\end{proposition}
\begin{proof}
Ellipticity needs no arithmetic at all: a diagram is elliptic if and only if each
of its connected components is one of the finite Coxeter types, which is a
condition on a labelled graph.

For Lann\'er membership, a diagram is Lann\'er if and only if it is connected, is
not elliptic, has all proper subdiagrams elliptic, and has $\det\Gr<0$; the last
clause is what separates it from parabolic, and given the third it also fixes the
inertia, because all proper subdiagrams elliptic makes every leading
$(k-1)$-minor positive, so by Jacobi's rule the sign of the determinant decides
between signature $(k,0)$, $(k-1,1)$ and a degenerate case. That sign is exact in
both cases that arise. On three nodes,
\[
  \sig\det\Gr \;=\; \sig\Bigl(\tfrac1p+\tfrac1q+\tfrac1r-1\Bigr),
\]
with a non-edge read as label $2$: the classical trichotomy for the triangle group,
and \emph{integer} arithmetic. A wild edge is not a label but an unknown $m_e\ge7$,
and it is covered too: the sum is evaluated at $m_e=7$ and is strictly decreasing in
$m_e$, so once it is below $1$ there it stays below $1$ for every $m_e\ge7$, and it
is below $1$ at $m_e=7$ for every connected three-node diagram, since two labels
$\ge3$ already give $1/p+1/q\le2/3$ and a label $2$ forces the other two $\ge3$. The
verdict is therefore constant on $m_e\ge7$, which is Lemma~\ref{lem:wildcard}(b). On four or five nodes no
label can exceed $5$ --- by Lemma~\ref{lem:wildcard}(c) a Lann\'er diagram of order
$\ge4$ has no edge of label $\ge6$ --- so the entries $-\cos(\pi/m)$ lie in
$K=\QQ(\sqrt2,\sqrt3,\sqrt5)$ and the determinant is computed exactly there. In
particular $\cos(\pi/7)$, which is cubic over $\QQ$, never enters: on three nodes
the integer criterion covers it, and on more nodes it is excluded outright.
\end{proof}

The Lann\'er gate therefore tests membership, where a floating-point gate would
naturally test the weaker necessary condition ``exactly one negative
eigenvalue''. That is sound, the subdiagram of a missing face being Lann\'er by
Proposition~\ref{prop:faces}; and what the strengthening adds is exactly the
no-parabolic-subdiagram clause of Proposition~\ref{prop:vinberg}, since the
condition it imposes beyond ``one negative eigenvalue'' is $\det\Gr\neq0$. So it
enforces a condition the polytope must satisfy anyway, rather than being merely
argued to be harmless. It also prunes harder, which is why every searched type here
is exhausted at the root of its label tree without any refinement
(\S\ref{sec:exhaust}).

\subsection{Symmetry breaking}
Let $\mathrm{Aut}(T)$ be the automorphism group of the facet--vertex incidence
structure of $T$, computed by VF2 subgraph isomorphism. A partial assignment is
discarded as soon as the vector of labels on its fixed prefix is provably not
lexicographically minimal among the images of that prefix under the
automorphisms stabilising the prefix's domain. Because the comparison is made on
the whole prefix vector rather than pair by pair, a discarded branch is
covered by the canonical branch in the same orbit, which is scheduled
separately; no solution orbit is lost. This is a throughput device only: on the two
most resistant $d=6$ types it turns prefixes that exhaust a $1800$~s enumeration
budget into sub-second certificates.

Because a failure of the device would be silent --- an orbit discarded whose
canonical member is never scheduled --- the verdicts are also checked with it
switched off. Seven of the $\NumRequired$ searched types are cheap enough to
enumerate without it, and all seven exhaust with unchanged verdicts. The case that
can fail in the damaging direction is the one where a polytope exists, since an
empty verdict is reproduced by a pruner that discards too much as readily as by a
correct one; that case is the subtree of the realizing type in which $\PP$ is
found. Run without pruning it enumerates $9{,}216$ complete labellings in place of
$952$, and returns the same verdict: exhausted, exactly one polytope, $\PP$.

\subsection{Deciding a labelling}
\label{sec:decide}
Given a complete ordinary labelling with no wildcard, the ultraparallel weights
are \emph{determined}, not searched: $\rank\Gr\le7$ forces every
$8\times 8$ minor of $\Gr$ to vanish, and for a real symmetric matrix it suffices
to impose this on the \emph{principal} minors, since the rank of such a matrix is
the largest order of a non-vanishing principal minor; and a minor whose index set contains
exactly one not-yet-pinned dashed edge is a quadratic in that single weight.
Solving it gives at most two roots $>1$; pinning cascades until all weights are
fixed, whereupon rank and signature are checked directly. When no
single-unknown minor is available, two minors over the same pair of unknowns are
combined by a Sylvester resultant, giving a quartic in one weight and a
back-substitution for the other.

Two points about the arithmetic of the screen are worth stating. On the single-unknown path the
minor is a quadratic in one variable, its coefficients are recovered by
interpolation at three points, and its roots are obtained in closed form. At a
\emph{tangential} solution that quadratic has a double root and its discriminant
is exactly zero, so in floating point the discriminant lands within a few units in
the last place of zero with a sign determined by the order of operations, and hence
by the numbering of the facets; a test of the form $\mathrm{disc}<0$ would then
refute a realizable labelling. We therefore compare each discriminant against a
relative tolerance $10^{-10}\cdot\max(b^2,|4ac|,1)$ and treat the near-zero case as
the double root it is. The error is one-sided in the safe direction: a spurious
candidate costs an exact certification that rejects it. Tangency is not a corner
case here---it is where the rank variety touches the constraint, which is to say
where rigid configurations live.

Note what Proposition~\ref{prop:plainexact} adds to this. Every rejection the
screen issues is required to carry an exact certificate, so a rejection the screen
gets wrong cannot pass unnoticed: it would be a labelling for which no certificate
exists, and it would be reported as such rather than silently dropped. The
discriminant tolerance above therefore affects how often the screen is useful, not
whether a verdict is right.

The pair-resultant path is weaker. Its quartic is recovered by least-squares
fitting through six sampled determinants and the degenerate,
identically-vanishing case is detected by a magnitude test; both operations are
scale-sensitive, so a refutation from that path is not reliable. We therefore trust
a refutation only when the pair-resultant path was not used, and otherwise treat
the outcome as inconclusive and fall back. This is strictly conservative, and by
Proposition~\ref{prop:cascade} it cannot alter any verdict in dimension $6$,
where that path is never taken.

If the cascade stalls, the code falls back to a numerical search: multistart
L-BFGS-B minimisation of the sum of squares of the three smallest singular
values of $\Gr$ over the box $x_e\in[1.001,1000]$, rejecting the labelling when
the attained value exceeds $10^{-6}$. This is the only bounded range of
ultraparallel weights anywhere in the cascade, and it is never reached.

\begin{proposition}[the cascade neither stalls nor takes the pair path]
\label{prop:cascade}
For every one of the $\NumTypes$ generated combinatorial types --- not merely the
$\NumRequired$ that reach the search:
\begin{enumerate}
\item[(i)] the dashed edges do not form a perfect matching, and
\item[(ii)] at every step of the cascade there is an $8\times8$ principal minor
      whose index set contains exactly one not-yet-pinned dashed edge.
\end{enumerate}
Consequently, in dimension $6$ the pair-resultant branch is never entered and the
cascade never runs out of pinnable dashed edges, so the bounded-box fallback is
structurally unreachable.
\end{proposition}
\begin{proof}
Both assertions are finite checks on the missing-face data of the $\NumTypes$
types, and neither depends on the labels or on the weight values: which $8\times8$
principal minors exist, which dashed edges each of their index sets contains, and
which dashed edges remain unpinned are all determined by the dashed-edge set alone.
The cascade's progress is therefore decidable statically, per type.

For (i), a perfect matching would need five pairwise-disjoint dashed edges
covering all ten facets. The dashed-edge counts of the $\NumTypes$ types range over
$2,\dots,6$, and enumerating the types with five or six shows none has five
pairwise-disjoint ones.

For (ii), writing $D$ for the dashed edges, such an index set exists as soon as some
$e\in D$ can be completed to an $8$-set meeting no other dashed edge in both
endpoints, which is decided by inspection of $D$; carrying out that decision for
each of the $\NumTypes$ types finds one in every case. What makes the completion
easy to exhibit is that $D$ is small: by Lemma~\ref{lem:degree} the dashed graph has
maximum degree at most $2$, so $|D|\le 10$. The check is nevertheless performed over
all $\NumTypes$ types rather than appealing to that bound.

Since a stall of the cascade means precisely the failure of (ii), and the
pair-resultant branch is entered only after such a failure, both conclusions
follow.
\end{proof}

The single-unknown path, which is the path dimension $6$ takes, is not covered by
that proposition; it is the path whose arithmetic is discussed above.

The cascade can also return ``inconclusive'' for value-dependent reasons rather
than structural ones, and none of these was observed to fire: branch counters over
every subtree of the $\NumRequired$ searched types
record no stall of any kind and $0$ calls into the
bounded-box fallback. Appendix~\ref{app:counters} gives the taxonomy and the
measurements, including the contrasting $d=4$ figures that show the instrumentation
does register the machinery where it is reachable.

The cascade evaluates its minors in double precision, so it serves as a screen and
no emptiness verdict rests on it. A labelling it does not accept is refuted
exactly.

\begin{proposition}
\label{prop:plainexact}
Each of the $\PlainRejected$ labellings without a wildcard that the cascade does
not accept is unrealizable, and this is certified exactly: for each there is an
$8\times8$ principal minor of $\Gr$ which is a nonzero element of
$K=\QQ(\sqrt2,\sqrt3,\sqrt5)$ determined by the ordinary labels alone, so
$\rank\Gr\ge d+2$ whatever the ultraparallel weights.
\end{proposition}

\begin{proof}
Realizability requires $\rank\Gr\le d+1$, hence the vanishing of every
$(d+2)\times(d+2)$ principal minor. Suppose the index set $S$ of such a minor
contains no dashed edge, that is, no pair $e\in D$ with both endpoints in $S$.
Then $\Gr[S,S]$ has no unknown entry: its off-diagonal entries are the ordinary
values $-\cos(\pi/m)$ with $m\in\{2,\dots,6\}$, all of which lie in $K$, so its
determinant is an element of $K$ and is computed exactly by rational arithmetic on
$8$-tuples in the basis $\{\sqrt{\prod P}:P\subseteq\{2,3,5\}\}$. If that element
is nonzero the labelling is refuted outright, with no weight variable, no interval
arithmetic and no tolerance entering the decision.

Such an index set exists here. These labellings all belong to the realizing type
(\S\ref{sec:wildvalid}), whose three dashed edges form a path on four facets; that
path has exactly three $2$-element vertex covers, so exactly three $8$-subsets of
the ten facets avoid every dashed edge. For each of the $\PlainRejected$
labellings at least one of those three determinants is nonzero, which is a finite
exact computation and is carried out for each. The one labelling for which all
three vanish is the one the search accepts, namely $\PP$.
\end{proof}

The certificate is therefore independent of the tolerances the cascade uses ---
discriminants near zero, root tests, vanishing thresholds --- because it forms none
of them; in particular the tangency comparison described above, on which a false
rejection would turn, cannot affect it. The residual dependence of an emptiness
verdict on floating point lies in the forward-checking prune, not in the screen;
see \S\ref{sec:assumptions}\ref{as:screen}.

\subsection{Wildcards}
\label{sec:wildcards}
A labelling containing $\star$ on a set $W$ of edges is not decided at $m=7$.
Each wild entry becomes a continuous unknown $c_e=\cos(\pi/m_e)\in[\cos(\pi/7),1)$
and the rank conditions are solved for the $c_e$ jointly with the dashed
weights. For each edge the feasible interval of $c_e$ is converted into an
integer window of admissible $m_e$, scanned up to $m_e=100$; every feasible
integer tuple is instantiated and passed to exact certification. If any window
reaches the scan limit the run raises a diagnostic flag and the type's verdict is
not counted as rigorous. \emph{The flag never fired};
observed windows closed by $m\approx30$.

The first step is a screen, and it disposes of a large part of the search: of the
$\NumAssignments$ complete labellings enumerated, $\NumWild$ carry at least one
wildcard, and none of them reaches the integer window scan. As a screen it is a
multistart quasi-Newton minimisation of the sum of squares of
the $n-d-1$ smallest singular values of $\Gr$, over $x_e\in[1.001,1000]$ for the
dashed weights and $c_e\in[\cos(\pi/7),1)$ for the free wild entries, from $24$
starts, rejecting when no start attains $10^{-6}$.

The failure of a local search over a bounded box is not a proof of infeasibility,
and no verdict rests on one: each of those $\NumWild$ labellings is refuted by an
exact certificate.

\begin{proposition}
\label{prop:wildexact}
Each of the $\NumWild$ wildcard-bearing labellings is infeasible, and this is
certified exactly: no local search, no bound on the ultraparallel weights, and no
floating-point tolerance enters the refutation.
\end{proposition}

\begin{proof}
Fix such a labelling, with dashed edges $D$ and wild edges $W$. Realizability
requires $\rank\Gr\le d+1$, hence the vanishing of \emph{every} $(d+2)\times(d+2)$
principal minor of $\Gr$, so exhibiting one principal minor that vanishes nowhere
on the domain refutes the labelling.

First compactify. Substituting $x_e = 1/t_e$ maps $x_e\in(1,\infty)$ to
$t_e\in(0,1)$. Each unknown occupies exactly the two symmetric positions of one
edge, so a principal minor is a polynomial of degree at most $2$ in each $x_e$ and
in each $c_e$; multiplying the minor on the index set $S$ by $t_e^{2}$ for each
$e\in D$ with both endpoints in $S$ therefore turns it into a polynomial in the
variables $(t_e)_{e\in D}$, $(c_e)_{e\in W}$. On the domain $t_e>0$, so this
changes no zero there; on the closed box it can only \emph{add} zeros, at
$t_e=0$, which makes a certificate harder to obtain and never makes one unsound.
The domain is contained in the closed rational box
\[
t_e\in[0,1]\quad(e\in D),\qquad c_e\in[c_7,1]\quad(e\in W),
\]
where $c_7$ is any rational lower bound for $\cos(\pi/7)$. Both enlargements ---
closing the box and lowering $\cos(\pi/7)$ to a rational --- only \emph{grow} the
domain, so emptiness on the box implies emptiness on the domain. The unbounded
direction has been removed rather than truncated, and no bound on $x_e$ is
assumed.

Next, the coefficients are exact. The ordinary entries $-\cos(\pi/m)$ for
$m\in\{2,\dots,6\}$ lie in $K=\QQ(\sqrt2,\sqrt3,\sqrt5)$, hence so do the
coefficients of every minor. For a minor with $r$ unknowns, which has degree at
most $2$ in each, those coefficients are recovered by tensor-product interpolation
from $3^{r}$ determinant evaluations over $K$ at integer nodes; $K$ is a field of
degree $8$ over $\QQ$ with the basis $\{\sqrt{\prod S}: S\subseteq\{2,3,5\}\}$, and
all of this is rational arithmetic on $8$-tuples.

Finally, emptiness on the box is certified by interval arithmetic. Enclosing each
coefficient by rationals and evaluating with \emph{outward}-rounded rational
intervals yields an enclosure of the polynomial's range over the box that is
guaranteed to be a superset of the true range; if that enclosure excludes $0$, the
box contains no zero, and the conclusion is a proof rather than an estimate. In
each of the $\NumWild$ cases a single $(d+2)$-principal minor settles the matter on
the root box, with no subdivision required. All but $\WildRootBox$ of them are
settled without interval arithmetic at all: the minor's index set contains neither a
dashed nor a wild edge, so the minor is a nonzero element of $K$ and the ordinary
labels alone force $\rank\Gr\ge d+2$. For the remaining $\WildRootBox$ the minor's
range over the whole box excludes $0$, in every case by a margin of at least
$\WildMinMargin$.
\end{proof}

Two things are worth drawing out. The certificate covers every integer $m_e\ge 7$
at once, because the wild entries enter only through
$c_e=\cos(\pi/m_e)\in[\cos(\pi/7),1)$, which is a \emph{compact} interval, and the
integer scan happens only after feasibility: infeasibility of the continuous
relaxation is already infeasibility for every admissible integer label, so no
unbounded quantifier over labels arises anywhere. And the certificate makes no use
of the box $[1.001,1000]$, so with Proposition~\ref{prop:cascade} no verdict reported
in this paper depends on a bounded range of ultraparallel weights, for labellings
with a wildcard as well as without.

A certificate of this kind can only ever refute; the property the $d=6$
conclusion needs is that it never refutes something realizable. That is
one-sidedness, and it is measured against ground truth on the $d=5$ census, on the
instances the screen itself saw.

Three counts have to be kept apart there. Of the $\DFiveWildCensus$
wildcard-bearing assignments in that census, the joint-feasibility step passes
$\DFiveWildScreenPass$; of those, $\DFiveWildSurvived$ become polytopes ---
exactly the $d=5$ polytopes carrying a label outside $\{2,\dots,6\}$, the nine with
a dihedral angle $\pi/10$ --- while the other
$\DFiveWildScreenPass-\DFiveWildSurvived$ are feasible for the \emph{continuous}
relaxation and are eliminated by the integer window scan, no admissible $m_e\ge7$
existing for them. Since the certificate refutes the continuous relaxation, it must
leave both groups undecided, so the statement to test is about all
$\DFiveWildScreenPass$ instances that survive the screen and not only the
$\DFiveWildSurvived$ that realize.

The certifier is run on the three types of that census in which every screen-passer
occurs: on all $\DFiveWildScreenPass$ instances that survive the screen anywhere in
the census, which is what the one-sided claim needs --- those are the only instances
a false refutation could damage --- together with $\DFiveWildJointKilled$ of their
joint-step rejections, sampled uniformly at a fixed seed. It refutes
$\DFiveWildPassRefuted$ of the survivors and $\DFiveWildCert$ of the
$\DFiveWildJointKilled$ rejections. So the refutation is one-sided in the safe
direction wherever the answer is known, and on the other side it is not weak: it
settles every rejection put to it. That also disposes of the concern that a blanket
rejector would produce the same $d=6$ output as a correct one.

The calibration exposes one limitation of the method: a continuous-relaxation
certificate cannot by itself close a labelling that is continuously feasible and
integrally infeasible, as those
$\DFiveWildScreenPass-\DFiveWildSurvived$ instances are. In dimension $6$ that
limitation is not reached, every one of the $\NumWild$ wildcard-bearing labellings
being refuted outright, but it is a feature of the method rather than an artefact
of these instances.

\subsection{Exact certification of an accepted labelling}
\label{sec:exactcert}
A labelling is accepted only after all of the following succeed.
\begin{enumerate}
\item Gauss--Newton refinement of the dashed weights at $100$ decimal digits on
      the overdetermined system of all $8\times8$ principal minors of $\Gr$
      meeting a dashed edge, with the ordinary entries evaluated as exact
      $-\cos(\pi/m)$ in \texttt{mpmath}; the final residual must be below
      $10^{-50}$.
\item Signature $(6,1)$ with exactly three zero eigenvalues at $100$ digits.
\item Local isolation: the Jacobian of the kernel conditions with respect to the
      dashed weights has full rank $k$. This rules out points on a
      positive-dimensional spurious component, on which the minor equations
      vanish identically and Gauss--Newton would converge to an arbitrary point.
\item No parabolic subdiagram (compactness).
\item An independent \texttt{CoxIter} \cite{Guglielmetti2017,GuglielmettiLMS2015}
      check certifying cocompactness and dimension $6$.
\end{enumerate}
The minimal polynomial of each weight is then recovered by integer relation
detection, giving exact algebraic data for every accepted weight.

\subsection{Exhaustion certificates}
\label{sec:exhaust}
The enumerator reports ``exhausted'' for a run only when every
branch was explored: no enumeration timeout, no cap on the number of
assignments, no break on the solve budget, and no wildcard deadline. A type
whose root run does not exhaust is split by fixing the first $j$ labels of the
deterministic enumeration order, producing six disjoint subtrees, recursively,
to depth at most $14$. Since the six labels partition the alphabet at each
position, the subtrees of a node partition its search space exactly, so a type's
verdict is rigorous precisely when its refinement tree is covered by exhausted
leaves. Every subtree result is checkpointed. In the event no type needs
refinement at all: every one of the $\NumRequired$ searched types is exhausted at
the root of its own label tree, so the refinement depth reached is $\MaxDepth$
against a cap of $14$, and the coverage recursion --- re-run independently over the
checkpoint file --- confirms all $\NumRequired$ types covered with no errors
recorded. The partitioning machinery is described because a type that did not
exhaust would need it, and because it is what the $d=4$ and $d=5$ censuses use.

\section{The classification theorem}
\label{sec:polytope}

\begin{theorem}
\label{thm:search}
Of the $\NumRequired$ combinatorial types that survive the restrictions of
\S\ref{sec:completeness}, exactly one---the type with $p=\RealizingP$ and
missing-face profile $\RealizingProfile$---admits a compact Coxeter structure, and
it admits exactly one up to isometry. The other $\NumRequiredLess$ are rigorous
zeros by exhaustion.
\end{theorem}

Together with Lemmas~\ref{lem:degree}, \ref{lem:mfsix}, \ref{lem:prism}
and~\ref{lem:facet}, this
gives Theorem~\ref{thm:main}.

The run screened $\NumAssignments$ labellings to
zero in $\CPUHours$ CPU-hours, across $\NumSubtrees$ certified subtree results ---
one per searched type, each exhausted at its root, so no type required refinement.

Two steps separate ``exactly one labelling is realizable'' from the uniqueness
claim in Theorem~\ref{thm:search}, and neither is covered by the refutation
certificates, so we report both.

First, a labelling need not determine its ultraparallel weights: each pinning step
solves a quadratic and admits up to two roots $>1$, so the weight tree of one
labelling may have several leaves and several could pass. For the accepted labelling
it has exactly $\WeightLeaves$ leaf, which is to say every pinning step had a single
admissible root, and that leaf is the one certified in \S\ref{sec:exactcert}. So the
labelling determines its Gram matrix outright.

Second, several labellings of a type can describe the same polytope, namely the
orbit of one under the combinatorial automorphism group $\mathrm{Aut}(T)$. Here
$|\mathrm{Aut}(T)|=\AutTypeOrder$ and the subgroup preserving the accepted labelling
has order $\AutStabOrder$, so the orbit consists of $\NumGramConfigs$ label
assignments --- exactly the $\NumGramConfigs$ the search returns when symmetry
breaking is switched off, and they give one Gram matrix up to simultaneous
permutation of rows and columns. A compact hyperbolic Coxeter polytope is determined
up to isometry of $\HH^6$ by its Gram matrix up to simultaneous permutation, so the
type realizes exactly one polytope up to isometry. The subgroup of order
$\AutStabOrder$ is the symmetry of the polytope's own labelled diagram.

The realizing polytope is not new: it is the one drawn as Figure~5 of
\cite{Burcroff2024}, attributed there to Bugaenko \cite{Bugaenko1984}, and we
verified that the polytope our search returns is isomorphic to that figure edge
for edge and label for label. We therefore refer to that figure rather than
redrawing it. Its Gram matrix, the automorphism computation behind ``exactly one
up to isometry'', and its position relative to the Felikson--Tumarkin class are
recorded with the data of \cite{Repo}, where the Gram matrix is also checked in
one command, independently of the search.

Table~\ref{tab:pertype} lists all $\NumSurvivors$ combinatorial types surviving
the facet filter of Lemma~\ref{lem:facet}, with the number $p$ of disjoint facet
pairs, the missing-face size profile, the verdict, the method by which it was
reached, the number of exhaustion certificates produced, the number of Coxeter
labellings enumerated, and the CPU time. Type identifiers refer to
the type list of \cite{Repo}. Every row is rigorous:
either a theorem or a fully covered refinement tree.

Two conventions in reading the table. First, ``labellings'' counts \emph{complete}
Coxeter labellings that reached the screen, so an entry of $0$ means
forward-checking pruned every branch before completion. It does not mean the type
was cheap: a type can be emptied by the prune and still cost hours, the cost being
the size of the tree explored before every branch died rather than the number of
branches that survived. Second, a bullet in the ``req.'' column marks the
$\NumRequired$ types that survive every restriction of \S\ref{sec:completeness} and
are therefore decided by search; the realizing type is among them. The remaining
rows are decided by a theorem, and carry no subtree or labelling counts.

{\small
\begin{longtable}{r r l c c l r r r}
\caption{The $\NumSurvivors$ surviving combinatorial types and their
verdicts.\label{tab:pertype}}\\
\toprule
type & $p$ & missing-face profile & req. & \#polytopes & method & subtrees &
labellings & CPU\,h\\
\midrule
\endfirsthead
\multicolumn{9}{l}{\small\emph{Table \ref{tab:pertype}, continued}}\\
\toprule
type & $p$ & missing-face profile & req. & \#polytopes & method & subtrees &
labellings & CPU\,h\\
\midrule
\endhead
\midrule
\multicolumn{9}{r}{\small\emph{continued on next page}}\\
\endfoot
\bottomrule
\endlastfoot
8 & 3 & $2^{3}3^{2}4^{2}5^{2}$ & $\bullet$ & 0 & search & 1 & 0 & 0.02 \\
12 & 3 & $2^{3}3^{2}4^{2}5^{2}$ & $\bullet$ & 0 & search & 1 & 0 & 0.01 \\
17 & 3 & $2^{3}3^{1}4^{4}5^{1}$ &  & 0 & Lem.\,\ref{lem:prism} & 0 & 0 & 0.00 \\
34 & 5 & $2^{5}3^{1}5^{1}$ &  & 0 & Lem.\,\ref{lem:degree} & 0 & 0 & 0.00 \\
36 & 4 & $2^{4}3^{1}4^{2}5^{1}$ &  & 0 & Lem.\,\ref{lem:degree} & 0 & 0 & 0.00 \\
38 & 4 & $2^{4}3^{2}5^{2}$ &  & 0 & Lem.\,\ref{lem:degree} & 0 & 0 & 0.00 \\
40 & 4 & $2^{4}3^{2}5^{2}$ &  & 0 & Lem.\,\ref{lem:degree} & 0 & 0 & 0.00 \\
51 & 4 & $2^{4}3^{1}4^{2}5^{1}$ &  & 0 & Lem.\,\ref{lem:degree} & 0 & 0 & 0.00 \\
55 & 5 & $2^{5}3^{1}5^{1}$ &  & 0 & Lem.\,\ref{lem:degree} & 0 & 0 & 0.00 \\
59 & 4 & $2^{4}3^{1}4^{2}5^{1}$ &  & 0 & Lem.\,\ref{lem:degree} & 0 & 0 & 0.00 \\
60 & 3 & $2^{3}3^{2}4^{2}5^{2}$ &  & 0 & Lem.\,\ref{lem:prism} & 0 & 0 & 0.00 \\
61 & 4 & $2^{4}3^{1}4^{2}5^{1}$ &  & 0 & Lem.\,\ref{lem:mfsix} & 0 & 0 & 0.00 \\
69 & 3 & $2^{3}3^{2}4^{2}5^{2}$ &  & 0 & Lem.\,\ref{lem:prism} & 0 & 0 & 0.00 \\
70 & 4 & $2^{4}3^{2}5^{2}$ &  & 0 & Lem.\,\ref{lem:degree} & 0 & 0 & 0.00 \\
92 & 5 & $2^{5}4^{2}$ &  & 0 & Lem.\,\ref{lem:degree} & 0 & 0 & 0.00 \\
103 & 5 & $2^{5}3^{1}5^{1}$ &  & 0 & Lem.\,\ref{lem:degree} & 0 & 0 & 0.00 \\
120 & 4 & $2^{4}3^{2}5^{2}$ &  & 0 & Lem.\,\ref{lem:degree} & 0 & 0 & 0.00 \\
127 & 3 & $2^{3}3^{3}5^{3}$ & $\bullet$ & 0 & search & 1 & 0 & 0.20 \\
132 & 5 & $2^{5}3^{1}5^{1}$ &  & 0 & Lem.\,\ref{lem:degree} & 0 & 0 & 0.00 \\
140 & 4 & $2^{4}3^{1}4^{2}5^{1}$ &  & 0 & Lem.\,\ref{lem:degree} & 0 & 0 & 0.00 \\
154 & 4 & $2^{4}3^{1}4^{2}5^{1}$ &  & 0 & Lem.\,\ref{lem:degree} & 0 & 0 & 0.00 \\
159 & 6 & $2^{6}$ &  & 0 & Lem.\,\ref{lem:degree} & 0 & 0 & 0.00 \\
162 & 3 & $2^{3}3^{1}4^{4}5^{1}$ &  & 0 & Lem.\,\ref{lem:prism} & 0 & 0 & 0.00 \\
168 & 3 & $2^{3}3^{1}4^{4}5^{1}$ &  & 0 & Lem.\,\ref{lem:prism} & 0 & 0 & 0.00 \\
173 & 3 & $2^{3}3^{3}5^{3}$ & $\bullet$ & 0 & search & 1 & 0 & 0.32 \\
206 & 4 & $2^{4}3^{2}5^{2}$ &  & 0 & Lem.\,\ref{lem:mfsix} & 0 & 0 & 0.00 \\
214 & 3 & $2^{3}3^{2}4^{2}5^{2}$ &  & 0 & Lem.\,\ref{lem:prism} & 0 & 0 & 0.00 \\
218 & 5 & $2^{5}4^{2}$ &  & 0 & Lem.\,\ref{lem:degree} & 0 & 0 & 0.00 \\
220 & 3 & $2^{3}3^{2}4^{2}5^{2}$ & $\bullet$ & 0 & search & 1 & 0 & 0.02 \\
229 & 4 & $2^{4}3^{2}5^{2}$ &  & 0 & Lem.\,\ref{lem:degree} & 0 & 0 & 0.00 \\
234 & 4 & $2^{4}3^{2}5^{2}$ &  & 0 & Lem.\,\ref{lem:degree} & 0 & 0 & 0.00 \\
239 & 4 & $2^{4}4^{4}$ &  & 0 & Lem.\,\ref{lem:degree} & 0 & 0 & 0.00 \\
255 & 3 & $2^{3}3^{2}4^{2}5^{2}$ & $\bullet$ & 0 & search & 1 & 0 & 0.00 \\
265 & 4 & $2^{4}3^{1}4^{2}5^{1}$ &  & 0 & Lem.\,\ref{lem:mfsix} & 0 & 0 & 0.00 \\
273 & 3 & $2^{3}3^{1}4^{4}5^{1}$ &  & 0 & Lem.\,\ref{lem:prism} & 0 & 0 & 0.00 \\
284 & 4 & $2^{4}3^{2}5^{2}$ &  & 0 & Lem.\,\ref{lem:mfsix} & 0 & 0 & 0.00 \\
286 & 5 & $2^{5}3^{1}5^{1}$ &  & 0 & Lem.\,\ref{lem:degree} & 0 & 0 & 0.00 \\
287 & 3 & $2^{3}3^{3}5^{3}$ & $\bullet$ & 0 & search & 1 & 0 & 0.04 \\
295 & 3 & $2^{3}3^{2}4^{2}5^{2}$ &  & 0 & Lem.\,\ref{lem:prism} & 0 & 0 & 0.00 \\
297 & 5 & $2^{5}4^{2}$ &  & 0 & Lem.\,\ref{lem:degree} & 0 & 0 & 0.00 \\
308 & 4 & $2^{4}4^{4}$ &  & 0 & Lem.\,\ref{lem:mfsix} & 0 & 0 & 0.00 \\
315 & 4 & $2^{4}3^{2}5^{2}$ &  & 0 & Lem.\,\ref{lem:mfsix} & 0 & 0 & 0.00 \\
317 & 3 & $2^{3}3^{2}4^{2}5^{2}$ &  & 0 & Lem.\,\ref{lem:prism} & 0 & 0 & 0.00 \\
320 & 3 & $2^{3}3^{3}5^{3}$ &  & 0 & Lem.\,\ref{lem:prism} & 0 & 0 & 0.00 \\
329 & 6 & $2^{6}$ &  & 0 & Lem.\,\ref{lem:degree} & 0 & 0 & 0.00 \\
332 & 3 & $2^{3}3^{2}4^{2}5^{2}$ & $\bullet$ & 0 & search & 1 & 0 & 0.01 \\
344 & 3 & $2^{3}4^{6}$ &  & 0 & Lem.\,\ref{lem:prism} & 0 & 0 & 0.00 \\
352 & 3 & $2^{3}3^{3}5^{3}$ & $\bullet$ & 0 & search & 1 & 0 & 0.00 \\
354 & 4 & $2^{4}4^{4}$ &  & 0 & Lem.\,\ref{lem:degree} & 0 & 0 & 0.00 \\
356 & 3 & $2^{3}3^{2}4^{2}5^{2}$ &  & 0 & Lem.\,\ref{lem:prism} & 0 & 0 & 0.00 \\
360 & 3 & $2^{3}3^{1}4^{4}5^{1}$ &  & 0 & Lem.\,\ref{lem:prism} & 0 & 0 & 0.00 \\
378 & 4 & $2^{4}4^{4}$ &  & 0 & Lem.\,\ref{lem:mfsix} & 0 & 0 & 0.00 \\
379 & 3 & $2^{3}3^{3}5^{3}$ & $\bullet$ & \textbf{1} & search & 1 & 952 & 0.27 \\
382 & 3 & $2^{3}3^{3}5^{3}$ & $\bullet$ & 0 & search & 1 & 0 & 0.05 \\

\end{longtable}
}

\section{Verification and reproducibility}
\label{sec:verification}

\subsection{The $d=5$ census, reproduced through the identical code path}
\label{sec:d5valid}

\emph{This is the anchor of record for the correctness of the code.} It is the
one validation that exercises the same functions, the same flags and the same
acceptance test as the $d=6$ classification, on a published census of known
value; where the two disagree, the code is wrong. The $d=4$ comparison of
\S\ref{sec:d4valid} plays a different and weaker role, for reasons given there.
It is also the strongest evidence we can offer that the search does not lose
polytopes. All $\DFiveTypes$ combinatorial types of simple $5$-polytopes with
$9$ facets and $p\ge2$ were run through the same per-type driver, in the same wildcard mode and with the
same acceptance test, and the outcome is Ma--Zheng's census exactly:
\[
\DFiveRealizing,\qquad\text{total } \DFiveTotal ,
\]
with every one of the $\DFiveTypes$ types reporting exhausted and no
window reaching the scan limit. The census is demanding in exactly the ways this
pipeline can fail: it contains polytopes with high-degree algebraic
weights, prism-glued families, and dihedral angles $\pi/10$---and the $\pi/10$
angles are \emph{discovered} by the wildcard range analysis rather than assumed
in any alphabet. Reproduce with the $d=5$ validation driver.

The validation was then repeated with the low-weight caps of
Lemma~\ref{lem:burcroff55b} additionally enabled---the one solver flag the $d=6$
run sets that the baseline $d=5$ validation did not exercise---and reproduces
the census again: $\DFiveTotal$ polytopes, the same six realizing types with the
same per-type counts, all $\DFiveTypes$ types exhausted, no
window reaching the scan limit, in $\DFiveCPUHours$ CPU-hours. The $d=5$ anchor therefore
exercises exactly the flag combination used for the $d=6$ classification.


The implementation of the facet filter is validated in the same dimension. Its
$d=5$ analogue uses the completed $d=4$ classification: a facet of a $5$-polytope
with $9$ facets meeting all eight others is a compact Coxeter $4$-polytope with $8$
facets, and $14$ of those $30$ combinatorial types realize one, between them
exhibiting ten distinct missing-face profiles. Requiring every such facet to carry
one of those ten profiles discards $48$ of the $109$ $d=5$ types and \emph{none} of
the six that realize a polytope. It also retains $55$ types that realize nothing,
which is the one-sidedness asserted above, here measured rather than assumed. The
filter is markedly weaker in dimension $5$ than in dimension $6$---it removes $44\%$
of the types rather than $86\%$---because the classification it consults is
correspondingly smaller.

\subsection{The $d=4$ census, reproduced exactly}
\label{sec:d4valid}

The same code path reproduces the second completed census in this family. Our
implementation returns all $\DFourTarget$ compact hyperbolic Coxeter
$4$-polytopes with $8$ facets: the $30$ combinatorial types are reproduced
exactly, all $30$ are certified exhausted by the coverage recursion of
\S\ref{sec:exhaust} over $90$ subtree results and $17.4$ CPU-hours, and the $14$
realizing types agree \emph{type by type} with \cite[Appendix~A]{Burcroff2024},
whose $34$ rows---the $14$ realizing types together with $20$ realizing
none---sum to $\DFourTarget$ in agreement with her abstract and with
\cite[Thm.~1.1]{MaZheng2024}. The types are matched to her $G$-types by exact
missing-face hypergraph isomorphism, with no discrepancy.

Dimension $4$ is a second anchor alongside \S\ref{sec:d5valid}, and we are
explicit about what it does and does not cover. It exercises the pair-resultant
branch of the cascade heavily---$132{,}571$ uses---which
Proposition~\ref{prop:cascade} shows is never entered in dimension $6$, so that
part of the evidence is not transferable. It also contains a type with no
dimension-$6$ analogue at all: $G_4$ is the $3$-free type, the $4$-cube, and by
Lemma~\ref{lem:esselmann} a $3$-free compact Coxeter $6$-polytope needs at least
$12$ facets. Burcroff attributes the twelve polytopes of that type to Jacquemet
and Tschantz \cite{JacquemetTschantz2018}.

Two features of Ma--Zheng's published intermediate data corroborate assumptions
made here rather than results. Their candidate label vectors use exactly the
alphabet $\{2,3,4,5,6,7\}$ with $7$ standing for ``any $m\ge7$'', an independent
implementation of the truncation of Proposition~\ref{prop:mz35} and of our
wildcard convention. And every polytope we find in the type they call $P_{8,17}$
appears in their candidate list for it, each once up to that type's order-$8$
automorphism group, so the two enumerations agree where they overlap.

\subsection{The exact refutation certificates, and what they were tested against}
\label{sec:wildvalid}
Propositions~\ref{prop:plainexact} and~\ref{prop:wildexact} are computations, so
they need the same treatment as the rest of the pipeline: the instances they refute
must be the instances the search screened, and their arithmetic must be checked
against something independent.

The instances are not re-derived. The classification driver writes every labelling
that reaches a screen to an instance sink --- its ordinary labels, dashed edges and
wild edges, together with the verdict it was rejected on: the numerical residual
for a wildcard-bearing labelling, the cascade's decision for the others --- and the
certifiers read that dump, so the labellings certified are the labellings screened.
The dump agrees with the classification's own counters
($\NumAssignments$ labellings, $\NumWild$ of them wildcard-bearing and $\NumPlain$
not, one realizer, exhausted), and carries a completion marker recording its record
count, so a partial dump cannot be mistaken for a complete one. In dimension $6$
all $\NumAssignments$ complete labellings arise inside a single subtree of the
realizing type; the other subtrees are emptied by forward checking before any
complete labelling is reached.

Four checks bound on the certifiers themselves.
\begin{enumerate}
\item \emph{The field arithmetic.} Multiplication, inversion by conjugates and
      determinants over $K$ are checked against independent computations, and each
      tabulated entry $-\cos(\pi/m)$ against its floating-point value.
\item \emph{The minor polynomials.} For randomly generated labellings of the same
      shape, every interpolated minor is compared against a NumPy determinant of
      the corresponding submatrix at random points of the domain --- the check that
      the polynomial being certified really is the minor. A wrong polynomial is the
      one failure mode that would refute realizable labellings silently.
\item \emph{The enclosures.} Sampled values of each minor must lie inside the
      computed range, and each recorded certificate in the artifact is
      re-checked against a floating-point recomputation of the minor it names.
\item \emph{$\PP$ itself.} Run on the realizing labelling, with its three dashed
      weights left unknown, the certifier must fail to refute it, and does: the
      branch and bound localises rather than excludes. Sharper, the three minors
      that carry every certificate of Proposition~\ref{prop:plainexact} --- the ones
      whose index sets avoid all three dashed edges --- vanish \emph{identically} in
      $K$ for this labelling, so no escalation of the search could refute it either.
      Over the whole search the same property is a census rather than a spot check:
      of the $\NumAssignments$ labellings dumped, exactly one is left unrefuted, and
      it is the one the search accepts.
\end{enumerate}
The calibration against the $d=5$ census is reported in \S\ref{sec:wildcards}:
$\DFiveWildCert$ of $\DFiveWildJointKilled$ sampled joint-step rejections
certified, and none of the $\DFiveWildScreenPass$ instances that survive that step
refuted.

\subsection{Independent checks on the polytope}
\label{sec:checks}
The Gram matrix of the realizing type was verified three ways: (i) exact inertia
over $\QQ(\sqrt2,\sqrt5)$ by symmetric congruence elimination, with the sign of
every pivot decided exactly and no floating point anywhere;
(ii) exact rank and the relation $w_{78}=w_{67}^2-1$ in SymPy; (iii)
\texttt{CoxIter}, which certifies cocompactness, dimension $6$, the $f$-vector
and the covolume. A single command, the verification script,
performs (i)--(iii).

\subsection{Provenance of the software}
\label{sec:provenance}

Most of the software on which this paper rests was written by Claude, a large
language model developed by Anthropic, working interactively under the author's
direction. This includes the order-type and Gale-diagram generator, the
combinatorial filters, the label enumerator and its forward checking, the
screening cascade, the exact certification code, the drivers, and the
verification scripts of \S\ref{sec:verification}. The author set the objectives,
chose the mathematical strategy, reviewed the code and the output, and is
responsible for the contents of this paper.

We disclose this because it bears on how the results should be read, and the
right response is not to ask the reader to trust the code's provenance but to
make the code unnecessary to trust. Every reduction from the full search space to
the final $\NumRequired$ decided types is a theorem, proved here or cited by
number (\S\ref{sec:completeness}); no step of that chain depends on the
implementation being correct, only on the implementation doing what the theorem
permits. Every accepted polytope is certified by exact arithmetic and then
re-checked by \texttt{CoxIter}, an independent and widely used program that
shares no code with ours. The emptiness verdicts do not need the numerical parts of
the code trusted either, and it is worth being precise about the division of
labour: the forward-checking prune, which on its own empties ten of the
$\NumRequired$ searched types, decides two predicates that carry no tolerance
(Proposition~\ref{prop:exactgates}); and each of the $\NumAssignments$ labellings
that reach a screen --- all of them inside the eleventh type --- carries an exact
refutation certificate (Propositions~\ref{prop:plainexact}
and~\ref{prop:wildexact}), for $\PlainRejected$ of them the non-vanishing of a
single determinant over $\QQ(\sqrt2,\sqrt3,\sqrt5)$, checkable independently of
this software.
The pipeline reproduces the published $d=5$ census
exactly through the identical code path (\S\ref{sec:d5valid}). The final Gram
matrix can be verified in one command with no reference to the search at all
(\S\ref{sec:checks}).

Readers inclined to weigh provenance heavily should weigh \S\ref{sec:d5valid},
\S\ref{sec:d4valid} and \S\ref{sec:assumptions} more heavily still.

\subsection{Data and code}
\label{sec:datacode}
Everything needed to check the results of this paper is available at \cite{Repo}:
the pipeline, the $\NumTypes$ combinatorial types, the $\NumSurvivors$-type
survivor list together with the facet filter that produces it, the $\NumSubtrees$
subtree certificates and per-type verdicts, the realizer records for the realizing
type, every labelling that reached a screen together with its exact refutation
certificate (\S\ref{sec:wildvalid}), the $d=5$ and $d=4$ census outputs,
and the verification scripts. Its \texttt{README} gives the commands, the expected
output of each, and the correspondence between the sections of this paper and the
code that produced them.

Two of those scripts are worth distinguishing from the rest, because neither
imports the pipeline: one re-verifies the Gram matrix of the realizing type from
the certificates alone, and one re-derives every number quoted here, recomputing
the coverage recursion of \S\ref{sec:exhaust} from the subtree certificates and
failing if any searched type is not closed, if any verdict rests on a truncation
at the refinement depth cap, on a wildcard window that reached its scan edge or on
a per-assignment deadline, if any labelling that reached a screen lacks an exact
refutation certificate, or if any labelling that survived a screen carries one.

\section{Assumptions and limitations}
\label{sec:assumptions}

We list, without softening, everything Theorem~\ref{thm:main} depends on beyond
the arguments given above.

\begin{enumerate}
\item\label{as:lit} \emph{Published theorems.} Lemma~\ref{lem:p2}
  (\cite{FeliksonTumarkin2009}), Lemma~\ref{lem:esselmann}
  (\cite[Lemma~6.7]{Esselmann1994}), Theorem~\ref{thm:burcroff35}
  (\cite[Thm.~3.5]{Burcroff2024}), Lemma~\ref{lem:burcroff55b}
  (\cite[Lemma~5.5(b)]{Burcroff2024}), Proposition~\ref{prop:mz35}
  (\cite[Prop.~3.5]{MaZheng2024}), Lann\'er's classification, and Vinberg's
  theory. Lemma~\ref{lem:prism}, which removes $\NumPrism$ types, rests in
  addition on the \emph{completeness} of the classification of compact hyperbolic
  Coxeter $d$-polytopes with $d+2$ facets --- Kaplinskaja's simplicial prisms
  \cite{Kaplinskaja1974} together with Esselmann's products of two simplices
  \cite{Esselmann1996}, the latter occurring only for $d=4$. Concretely it needs
  that no compact Coxeter $\Delta^2\times\Delta^3$ exists in $\HH^5$. That is a
  citation and not a consequence of Lann\'er's classification, so we list it
  separately.

\item\label{as:facet} \emph{The facet filter rests entirely on
  \cite{MaZheng2023}.} Lemma~\ref{lem:facet} removes $333$ of $\NumTypes$
  types, and it is exactly as reliable as Ma--Zheng's assertion that precisely
  six simple $5$-polytopes with $9$ facets admit a compact Coxeter structure.
  That classification is corroborated by Burcroff's independent count of the
  same $51$ polytopes and by our own reproduction of the census
  (\S\ref{sec:d5valid}), which independently finds the same six realizing
  types. If a seventh realizing $5$-polytope with a fifth missing-face profile
  were ever found, the filter would have to be rerun. We have not re-derived
  Ma--Zheng's $d=5$ classification from scratch by an independent method.

\item\label{as:db} \emph{Order-type database.} We rely on
  \cite{AichholzerAurenhammerKrasser2002,OrderTypeDB} being a complete list of
  realizable order types on $10$ points (\S\ref{sec:aak}).

\item\label{as:screen} \emph{Emptiness verdicts.} Acceptances are exact
  (\S\ref{sec:search}) and independently checked by \texttt{CoxIter}. Rejections
  come from two mechanisms, and both are exact.

  \emph{Forward checking.} Most labellings never reach a screen: they are pruned
  during enumeration by the ellipticity condition on a face's subdiagram or the
  Lann\'er condition on a missing face's. Only a false \emph{rejection} would be
  unsound, and neither test can produce one, because neither is numerical.
  Ellipticity is a condition on a labelled graph and needs no arithmetic; Lann\'er
  membership reduces to a determinant sign that is an integer comparison on three
  nodes and an exact computation in $K=\QQ(\sqrt2,\sqrt3,\sqrt5)$ on four or five,
  the higher-degree case $\cos(\pi/7)$ being excluded outright by
  Lemma~\ref{lem:wildcard}(c) (Proposition~\ref{prop:exactgates}). The search
  decides both this way, so the prune carries no tolerance. As corroboration the
  exact tests were also compared against the floating-point tests they replace:
  they agree on all $216$ label tuples on $3$ nodes and all $46{,}656$ on $4$
  nodes, exhaustively, and the exact Lann\'er test implies the weaker
  one-negative-eigenvalue condition in every one of those cases.

  \emph{The cascade and the wildcard screen.} Both are evaluated in double
  precision --- the cascade with root and vanishing thresholds of $10^{-7}$ to
  $10^{-12}$ and a leaf test at $10^{-7}$ (singular value) and $10^{-6}$
  (signature), the wildcard step with a multistart floor at $10^{-6}$ --- so a
  rejection from either, taken on its own, could be a false negative caused by an
  ill-conditioned minor sitting near a tolerance. Neither is taken on its own.
  Every labelling that reaches a screen and is not accepted carries an exact
  refutation: Proposition~\ref{prop:wildexact} for the $\NumWild$ wildcard-bearing
  labellings and Proposition~\ref{prop:plainexact} for the $\PlainRejected$ others.
  Those two counts together with the single acceptance are the $\NumAssignments$
  labellings that reach a screen, and the accepted one is refuted by neither. Those
  certificates are exact arithmetic in $K$ with outward-rounded rational intervals,
  so no tolerance of the screens enters them, and in particular neither the
  pair-resultant branch, which is provably never reached at $d=6$
  (Proposition~\ref{prop:cascade}), nor the tangential discriminant compared
  against a relative tolerance (\S\ref{sec:decide}) bears on any verdict: the
  certificates form neither quantity.

  So no emptiness verdict reported here depends on floating-point arithmetic. What
  the screens' tolerances still affect is how much work the exact steps are left to
  do: a screen that wrongly rejected a labelling would produce a rejection for
  which no certificate exists, which is reported rather than silently accepted
  (\S\ref{sec:decide}), and a screen that wrongly accepted one costs an exact
  certification that rejects it.

  \item\label{as:wild} \emph{Wildcard resolution.} Of the $\NumAssignments$
  labellings enumerated, $\NumWild$ carry a wildcard, and each is refuted by
  Proposition~\ref{prop:wildexact}: a single principal minor per labelling, on a
  compactified box carrying no bound on the weights, with no floating-point
  tolerance in the decision. The screen that precedes it --- a multistart local
  minimisation over $x_e\in[1.001,1000]$ and $c_e\in[\cos(\pi/7),1)$, rejecting
  when no start reaches $10^{-6}$ --- would rest on two things the rest of the
  paper avoids, the failure of a local search and a bounded range of dashed
  weights, and no verdict rests on it. What we do rely on is that the
  certifier's arithmetic is what it claims to be: the exact field arithmetic over
  $K$, the interpolation of each minor, and the outward rounding of every
  interval operation. Those are tested against independent computations
  (\S\ref{sec:wildvalid}), including that the certifier does not refute $\PP$
  itself, and calibrated against the $d=5$ census, where the answer is known. That
  calibration is exhaustive on the side that decides one-sidedness --- all
  $\DFiveWildScreenPass$ instances that survive the screen anywhere in the census,
  none of them refuted --- and sampled on the other, $\DFiveWildCert$ of
  $\DFiveWildJointKilled$ rejections certified.
  The integer scan bounded at $m\le100$, guarded by a
  diagnostic flag raised whenever a window reaches that limit, is a separate matter
  and does not arise at $d=6$: no wildcard-bearing labelling reaches it.

\item\label{as:tools} \emph{Tool correctness.} We rely on \texttt{CoxIter} for
  the final compactness and dimension checks, and on NumPy/SciPy/SymPy/mpmath.

\item\label{as:aut} \emph{Automorphism pruning.} The orbit symmetry-breaking of
  \S\ref{sec:search} is argued sound there, and its verdicts are reproduced with
  the device switched off on the seven searched types that can be enumerated
  without it and --- the case that matters --- on the subtree in which $\PP$ is
  found, where without pruning ten times as many labellings are enumerated and the
  same single polytope comes out. Four searched types cannot be enumerated without
  it, which is why it exists; for those the argument is the lex-leader one and not
  a measurement. Its failure mode would be silent, so we name it explicitly.

\end{enumerate}

\appendix

\section{The six realizing $5$-polytopes with $9$ facets}
\label{app:d5types}

Lemma~\ref{lem:facet} filters a $d=6$ type by requiring every facet that meets all
nine others to have the missing-face profile of one of the six combinatorial types
of simple $5$-polytopes with $9$ facets that admit a compact Coxeter structure
\cite[\S6]{MaZheng2023}. Those six, and hence the four profiles the filter tests
against, are recorded here so that the filter can be checked without reconstructing
them.

Facets are numbered $0,\dots,8$; each entry of the last column is a missing face,
written as the concatenation of its facets. Numbering follows
\cite{MaZheng2023}, whose list of all $322$ simple $5$-polytopes with $9$ facets we
matched to ours by exact hypergraph isomorphism. Those $322$ are all simple
$5$-polytopes with $9$ facets; the $\DFiveTypes$ types of \S\ref{sec:genvalid} are
the subset with $p\ge2$, which by Lemma~\ref{lem:p2} is the only part a compact
Coxeter structure can occur in. Under the type numbering of our own list, used in
\S\ref{sec:d5valid}, the six correspond as
\[
P_{9,319}\!\to\!0,\quad P_{9,322}\!\to\!6,\quad P_{9,302}\!\to\!19,\quad
P_{9,313}\!\to\!29,\quad P_{9,312}\!\to\!5,\quad P_{9,284}\!\to\!63 ,
\]
which is the order in which \S\ref{sec:d5valid} reports their polytope counts.

\begin{center}
{\small
\begin{tabular}{l r l p{0.46\textwidth}}
\toprule
type & \#polytopes & profile & missing faces\\
\midrule
$P_{9,284}$ & 1 & $2^{3} 3^{3} 4^{3}$ & \texttt{57,\,58,\,68,\,015,\,018,\,467,\,0123,\,2346,\,2347} \\
$P_{9,302}$ & 6 & $2^{5} 3^{1} 4^{1} 5^{2}$ & \texttt{47,\,57,\,58,\,68,\,78,\,015,\,2346,\,01234,\,01236} \\
$P_{9,312}$ & 1 & $2^{6} 5^{3}$ & \texttt{34,\,46,\,47,\,56,\,67,\,78,\,01235,\,01238,\,01258} \\
$P_{9,313}$ & 3 & $2^{4} 3^{2} 4^{2} 5^{1}$ & \texttt{46,\,56,\,68,\,78,\,234,\,347,\,0125,\,0158,\,01237} \\
$P_{9,319}$ & 22 & $2^{6} 5^{3}$ & \texttt{46,\,47,\,57,\,58,\,67,\,68,\,01234,\,01235,\,01238} \\
$P_{9,322}$ & 18 & $2^{6} 5^{3}$ & \texttt{34,\,46,\,47,\,56,\,68,\,78,\,01235,\,01237,\,01258} \\
\bottomrule
\end{tabular}}
\end{center}

The four distinct profiles are $2^6 5^3$ (shared by $P_{9,312}$, $P_{9,319}$ and
$P_{9,322}$), $2^5 3\,4\,5^2$ ($P_{9,302}$), $2^4 3^2 4^2 5$ ($P_{9,313}$) and
$2^3 3^3 4^3$ ($P_{9,284}$), which are the four listed in Lemma~\ref{lem:facet}.

\section{Hardware and timings}
\label{app:timings}
All computations were run on a single Apple M1 (8 cores, 16\,GB RAM),
macOS~26.5.2, Python~3.10.0 with NumPy, SciPy, SymPy and mpmath, and
\texttt{CoxIter} built from source. The $d=6$ classification consumed
$\CPUHours$ CPU-hours across $\NumSubtrees$ subtree tasks, one per searched type
and each exhausted at its root, so the refinement depth reached is $\MaxDepth$
against a cap of $14$. Subtree tasks are independent and were run concurrently. The
$d=5$ revalidation consumed $\DFiveCPUHours$ CPU-hours over its $\DFiveTypes$
types and the $d=4$ census $17.4$ CPU-hours over its $30$. Per-type timings are
in Table~\ref{tab:pertype}.

One environmental factor is worth recording, because it dominates the timings and
is invisible in the algorithm. The determinants driving the screen are $8\times8$
at most, far too small for a threaded BLAS to repay its synchronisation cost, so
running seven worker processes each with a default thread pool oversubscribes the
machine badly. Pinning every worker to a single BLAS thread
(one thread per worker and its equivalents) left the arithmetic
bit-for-bit identical---verified by hashing the recovered weights of all $24$
orbit representatives of the $G_{11}$ polytopes at one and at eight threads---while
reducing the cost of one representative $d=4$ type from $963$ to $22$ seconds.

\section{Branch counters and stall analysis for the cascade}
\label{app:counters}

Proposition~\ref{prop:cascade} rules out the \emph{structural} stall of the cascade:
running out of pinnable dashed edges. The cascade can also return
``inconclusive'' for value-dependent reasons, and these are worth separating,
because two of the three are confined to the pair-resultant branch that
Proposition~\ref{prop:cascade} shows is never entered in dimension $6$. A minor
that vanishes identically in its unknown, and a search tree exceeding $256$ leaves,
can occur on the single-unknown path; an identically-vanishing resultant and the
absence of a pinnable pair cannot arise at all when no pair is ever pinned. So only
two routes to the bounded-box fallback are available in dimension $6$.

Neither was observed. The run records $\NumCascade$ cascade decisions, all but the
single acceptance being ordinary rejections: the pair-resultant branch is entered
$0$ times, no stall of any kind fires, and the bounded-box fallback runs $0$ times.
The same instrumentation on the $d=4$ census of record makes $\DFourCascade$
cascade decisions, of which $\DFourPair$ enter the pair-resultant branch and
$\DFourFallback$ call into the bounded-box fallback. Those two counts are equal
because they must be: a refutation from the pair-resultant branch is not trusted
(\S\ref{sec:decide}), so every use of that branch is downgraded to inconclusive
and falls through to the fallback. The machinery therefore does fire where it is
reachable, and its absence at $d=6$ is not an artefact of the counters.

The hypothesis of Proposition~\ref{prop:cascade} is not delicate, and it is worth
recording how far it holds. The cascade is forced onto the pair-resultant branch
exactly when no $8$-minor isolates a single dashed edge, which happens in dimension
$4$ for the type whose dashed edges form a perfect matching---the $4$-cube---and for
no generated $d=6$ type at all. Refutations from that branch are not trusted
(\S\ref{sec:decide}), which is why its unreachability in dimension $6$ is worth
establishing rather than assuming.

We stress the limit of these measurements. They cover every cascade decision of
the search, but they count only the branches the instrumentation was written to
count, so they establish that the cascade's bounded-box fallback is not load-bearing
for any $d=6$ verdict without being a proof of it, and they say nothing about the
arithmetic \emph{within} a branch that does fire. They also say nothing about the wildcard feasibility step of
\S\ref{sec:wildcards}, which searches a box of the same shape. That step needs no
mitigation from counters, because no verdict rests on it: the $\NumWild$
wildcard-bearing labellings are refuted by the exact certificates of
Proposition~\ref{prop:wildexact}, which use no box in the ultraparallel weights.

\bibliographystyle{amsplain}
\bibliography{references}

\end{document}
